\documentclass[11pt]{article}
\usepackage[margin=1in]{geometry}
\usepackage{amsmath,amssymb,amsthm}
\newtheorem{definition}{Definition}
\usepackage{booktabs}
\usepackage{graphicx}
\usepackage{array}
\usepackage{xcolor}
\usepackage{hyperref}
\usepackage{enumitem}
\usepackage{subcaption}
\usepackage{pifont}

% --- notation macros (carried from Chapter A v2 for consistency) ---
\newcommand{\R}{\mathbb{R}}
\newcommand{\Q}{\mathbf{Q}}
\newcommand{\X}{\mathbf{X}}
\newcommand{\uu}{\mathbf{u}}
\newcommand{\yy}{\mathbf{y}}
\newcommand{\Amat}{\mathbf{A}}
\newcommand{\W}{\mathbf{W}}
\newcommand{\bmu}{\boldsymbol{\mu}}
\newcommand{\bSigma}{\boldsymbol{\Sigma}}
\newcommand{\zz}{\mathbf{z}}
\newcommand{\cmark}{\ding{51}}
\newcommand{\xmark}{\ding{55}}

% \REDO{...}: a number that is REAL but was measured on a superseded checkpoint/protocol and must be
% re-run before submission. \PLAN{...}: an experiment that is designed and specified but NOT YET RUN --
% no number exists. Both render visibly so neither can survive to camera-ready unnoticed.
\newcommand{\REDO}[1]{\textcolor{red}{[REDO: #1]}}
\newcommand{\PLAN}[1]{\textcolor{blue}{[PLANNED, NOT YET RUN: #1]}}

\title{\textbf{What Should a Cold Interview Optimise?\\Objective Mismatch in Adaptive Preference
Elicitation}\\[4pt]\large Thesis Chapter B (v2) --- working draft}
\author{}
\date{August 2026}

\begin{document}
\maketitle

\begin{abstract}
\noindent A cold-start conversational recommender must choose what to ask. The literature's answer is
overwhelmingly \emph{adaptive}: condition each question on the answers so far. Measured on a certified
instrument against a static order built by maximising the metric we report, every published adaptive
selector we could replicate loses --- and this chapter's claim is that the reason is the objective,
not the answers. Our contributions are three. \textbf{(i) A replication null with a hard comparator.}
Against a prefix-optimal static order optimised on the same recommender it is scored on, the bandit
family (GLM-UCB and Thompson sampling over a matrix-factorisation
posterior~\cite{zhao2013interactive}), a myopic information-gain acquisition~\cite{ma2019eddi}, the
conversational-bandit line~\cite{zhang2020conucb,wang2023conlinucb,liu2025clisk} and Golbandi's
interview tree~\cite{golbandi2011adaptive} lose at every budget from one to sixteen questions on
$10{,}000$ held-out users, on four of the five recommenders we score them with, including the one each
was built for. The loss is not for want of branching: Thompson sampling gives all $10{,}000$ users a
distinct interview by the fourth question and still loses. On the fifth recommender the ordering
reverses, and we report why: which selector wins depends on which recommender one assumes.
\textbf{(ii) The diagnosis: an objective mismatch, measured rather than asserted.} Exactly one of the
twenty-one distinct selection rules in the battery optimises held-out ranking quality, and it is the
static order. Read against exact measured NDCG@10 gains at the $22$ states an adaptive tree actually
visits, the published criteria are close to uninformative: mean Spearman $+0.010$ (rating
variance~\cite{golbandi2011adaptive}), $-0.060$ (pairwise rank loss~\cite{hu2013interview}) and
$-0.104$ (the logistic loss that paper's experiments ran), with a user-weighted regret of $66$--$89\%$
of what a uniform random pick forgoes.
\textbf{(iii) The demonstration: change only the objective and adaptivity wins.} Substituting measured
held-out NDCG@10 for the rating-variance criterion in the field's own ternary interview-tree machinery
--- same tree code, same stopping rule, same charged behavioural answers, same recommender, same pool,
same cohorts, same budget --- beats the static order on held-out users \textbf{from four questions
onward on both axes} ($+0.01035$ full $[+0.00749,+0.01328]$ and $+0.00705$ tail at sixteen), while the
two published surrogates lose by three times that margin. The size qualifiers travel with the claim:
conditioning is worth about a tenth of the interview's own lift, it is \emph{smaller} than the gain
from changing the question bank with no conditioning at all ($+0.01430$ at equal budget), it is
exactly zero at one question by construction and behind on the tail at two, and $39.3\%$ of the test
cohort sits in a leaf that loses. Adaptivity pays; it is not large. Finally we report the reach
problem that conditioning does not touch: $\sim\!14\%$ of cold users answer nothing across sixteen
charged questions, and answering \emph{at all} is worth roughly $+0.25$ against $+0.048$ for the best
question \emph{choice}.
\end{abstract}

\section{Introduction}

\paragraph{The claim this chapter makes.} \emph{Published adaptive selectors lose to a well-chosen
fixed order at every budget we test --- not because cold, closed, charged answers fail to
differentiate users, but because none of those selectors optimises the quantity anyone reports.} The
first half is a replication result (Section~\ref{sec:act1}); the second half is a diagnosis we
measure rather than assert (Section~\ref{sec:act2}); and the two together yield a demonstration.
\textbf{Exactly one of the twenty-one distinct selection rules in this battery optimises held-out
ranking quality, and it is the static order.} The rest maximise information, minimise regret, or
reduce within-node rating variance. Swap only the split criterion in the field's own interview-tree
machinery~\cite{golbandi2011adaptive} for measured held-out NDCG@10 --- holding the tree code, the
ternary branch rule, the stopping rule, the answers, the recommender, the candidate mask, the pool,
both cohorts and the budget fixed --- and the sign of the conditioning gain flips: adaptivity beats
the strong static order on held-out users from four questions onward, on both axes
(Section~\ref{sec:act3}).

\textbf{The size qualifiers travel with that claim wherever it appears.} The conditioning gain at
sixteen questions is $+0.01035$ full and $+0.00705$ tail --- about a tenth of what the interview
itself buys over asking nothing, and \emph{smaller} than the gain from changing the question bank
with no conditioning at all ($+0.01430$ full at equal budget). At one question the tree \emph{is} the
static order by construction, at two it is behind on the tail, and $39.3\%$ of the test cohort lands
in a leaf that loses. Adaptivity pays. It is not large. And the conditioning it buys is not a
curriculum: the branches ask almost disjoint question sets, but a pre-registered test of whether a
branch's questions are \emph{about} the answer that created it returns a null against both its
controls (Section~\ref{sec:routing}), so no routing claim is made anywhere in this chapter.

\paragraph{What this chapter is not claiming.} Question choice is decisively important here:
the difference between a good and a bad static order is large (a fixed bank moves full NDCG@10 from
the popularity prior's $0.1626$ to $0.2303$ at sixteen questions), and nothing here licenses the
reading that question choice is unimportant. It is not claiming a field-wide adaptivity null: no
such published null exists, and our own sub-threshold measurements are results \emph{in} this
chapter, not background it can cite. It is emphatically \textbf{not} claiming that nobody ever
compared an adaptive selector against a metric-optimised static order --- Golbandi's own tree beat
GreedyExtend~\cite{golbandi2010bootstrapping}, a seed set built by greedily minimising the RMSE its
era reported, and beat it decisively. The claim is scoped to what happened afterwards: since the
field's reported metric moved to truncated, position-discounted top-$k$ ranking, no adaptive
selector has been tested against a static order built by maximising \emph{that} metric.
Section~\ref{sec:positioning} states this concession, and three others, in full.

\paragraph{Relationship to Chapter A.} Chapter A built and certified the \emph{instrument}: a
recommender that ingests an arbitrary-length, mixed-channel set of answers without per-user
retraining and without conceding full-profile accuracy (four requirements R1--R4). It explicitly
deferred policy: ``it does not claim a new elicitation policy (that is the companion method
chapter); it claims the recommender that makes policy measurement possible.'' This chapter is that
companion. Everything below runs on the certified pair --- tower \texttt{i26c ep31} (test full/tail
NDCG@10 $0.3496/0.2491$) and concept channel \texttt{cv3 ep08} ($464$k parameters) --- both frozen
and read-only throughout. No model is trained in this chapter.

\section{Setup}
\label{sec:setup}

\paragraph{Ruler.} The canonical strong-generalisation split of Liang et
al.~\cite{liang2018variational} applied to MovieLens-25M~\cite{harper2015movielens} restricted to
the project catalogue: $18{,}359$ train-vocabulary items, $140{,}768$ train / $10{,}000$ validation
/ $10{,}000$ test users, per-user 80/20 fold-in/target, seed $98765$. Primary metric: full
\emph{and} tail NDCG@10, the tail measured on a head-deleted candidate list under the
short-head definition of Cremonesi et al.~\cite{cremonesi2010performance}, scored with the
recommender's learned decoder bias. The full catalogue is ranked --- no negative sampling anywhere,
since sampled ranking metrics do not preserve the ordering of
methods~\cite{krichene2020sampled}. The interview contract is NO-FLOOR (a model that received no
answer is scored on what it actually predicts, not overwritten by the popularity row); the floored
arm is a robustness row, not the headline.

\paragraph{Answers are behavioural.} A question is answered from the user's own held-out ratings:
an item question by their half-star level if they rated it, a concept question by the signed
implicit-watch statistic over their profile. An asked question the user cannot answer \emph{burns
the budget} --- items fold as the tower's native unseen-exposure token, concepts as the concept
channel's refusal token. No answer is generated by a model: every answer \emph{value} is read from
the user's own recorded behaviour, and no language model appears anywhere in the measurement path.
This matters because a simulator that answers everything would erase the single most important
asymmetry in the results below. What is nonetheless \emph{inferred} rather than observed is
\emph{whether} a user could have answered at all --- we take having rated the item as the criterion
--- and Section~\ref{sec:limits} states the consequences.

\paragraph{Selection interface.} A \emph{selector} sees, per user, the questions asked so far and
the answers actually obtained, and returns the next question; a \emph{consumer} is the recommender
that folds the final answer set and ranks the catalogue. The two are independent, and we vary them
independently. Static orders are the degenerate case in which the selector ignores the answers.
Before any adaptive number was produced we verified that all six static arms driven through the
sequential loop reproduce the batch implementation \emph{exactly} --- list equality of asked items,
answered items and half-star levels, over $30$ (arm, budget) cells on all $10{,}000$ users.

\subsection{What the answer model can and cannot launder}
\label{sec:circularity}

The standing objection to any conversational-recommendation result is that the thing answering the
questions is a version of the thing being evaluated, so the interview is graded by its own
assumptions. It is the right objection to raise and we would raise it first, so we state exactly
where this study is exposed to it and where it is not.

\emph{Answer values are not modelled at all.} Every answer is read from the user's own recorded
behaviour --- a half-star level for an item, the signed implicit-watch statistic for a concept. No
simulator is fitted, no language model is called, and no answer is generated anywhere in the
measurement path. This is a stronger separation than the surrounding literature usually has, and it
is the reason the asymmetry Act 4 turns on survives at all: in a protocol where the simulator answers
from the target item's attribute set, every question is answerable by construction and the cost of an
unanswerable question is identically zero (Section~\ref{sec:positioning}).

\emph{Answerability is inferred, and it is load-bearing.} Whether a user \emph{could} have answered
is not recorded anywhere, so we infer it --- having rated the item, for an item question; a
structural coverage rule, for a concept. This is a modelling assumption, it is the one in this
chapter we would most like to remove, and because two of the mechanisms below depend on it we return
to it in Section~\ref{sec:limits} rather than leave it to a footnote.

\emph{The selector and the consumer share a frozen encoder, and we do not remove that.} The
representation used to choose a question is the representation used to score the recommendation.
What that shared representation cannot do is manufacture the effects this chapter reports, and the
reason is structural rather than rhetorical: every claim here is a \emph{difference} between two
arms that share the encoder equally. The comparator is not a weaker family --- it is the same greedy
search over the same bank, scored by the same recommender, differing only in whether the question
may condition on the answers so far. A representation that flatters the interview flatters both
arms and cancels in the difference. Three checks make that concrete rather than assumed. A tree grown
on a shuffled routing signal does not merely fail to help, it \emph{loses} --- $-0.04817$ full under
the min-leaf-respecting shuffle and $-0.03934$ unrestricted --- so the structure is doing work that
no shared representation supplies for free. Changing only the scalar inside the split, one line with
everything else held, flips the sign of the conditioning gain (Section~\ref{sec:criterion}). And the
effect survives rebuilding the comparator until it is strictly harder, which cost us a separated
interval we had already printed in our own favour.

What none of this buys is a claim about live users. A behavioural answer is a real answer the user
once gave, not the answer they would give in a conversation, and the gap between those two is not
measured here by anyone, us included.

\section{Act 1: the frontier under standard practice}
\label{sec:act1}

\subsection{The comparator is the hard one}

An adaptivity null is only interesting if the static bar is genuinely strong. Ours is the greedy
prefix-optimal sequence over a bank of $500$ popularity-ranked items and $500$ answerability-ranked
concepts, computed by lazy greedy maximisation of held-out NDCG@10 \emph{on the same recommender it
is scored on}. Chapter A establishes that greedy is prefix-optimal here, which is what makes
``strongest static'' a computable object rather than an aspiration. \textbf{Every order in this
chapter --- global, group-conditioned, single-channel and combined --- maximises \emph{full}
NDCG@10; tail NDCG@10 is reported throughout as a secondary outcome and never enters a search.}
The search ran on the
$10{,}000$ validation users only; the test cohort was not constructed until all sequences were
frozen. On those validation users the empty-interview row reproduces the tower's certified cold
number to five decimal places ($0.12987$ against a recorded $0.12987$), which is how we know the
harness is measuring the certified system. That bridge is computed under the recommender's own
certification pool; the interview tables that follow use the protocol pool defined below, under
which the same empty interview on the test cohort reads $0.1586/0.0282$ --- the intercept of
Table~\ref{tab:greedy}.

\begin{table}[t]
\centering
\caption{Greedy prefix-optimal static sequences on the certified stack, test cohort
($10{,}000$ users), full/tail NDCG@10, protocol candidate pool. Intercept (no questions):
$0.1586/0.0282$. All three sequences are selected and scored under the same candidate mask. At
$q{=}1$ and $q{=}2$ the combined order \emph{is} the concepts order --- its first two picks are the
same two concepts --- so those cells are equal by nesting, not by adjudication.}
\label{tab:greedy}
\footnotesize
\setlength{\tabcolsep}{2pt}
\begin{tabular}{llccccc c}
\toprule
arm & emergent order & $q{=}1$ & $q{=}2$ & $q{=}4$ & $q{=}8$ & $q{=}16$ & answered@16 \\
\midrule
items-only & \texttt{iiiiiiiiiiiiiiii} & .1760/.0434 & .1884/.0563 & .2038/.0677 & .2235/.0860 & .2447/.1017 & 3.11 i \\
concepts-only & \texttt{cccccccccccccccc} & \textbf{.1808}/\textbf{.0537} & \textbf{.1962}/\textbf{.0701} & .2141/\textbf{.0858} & .2347/\textbf{.1074} & .2447/.1161 & 12.36 c \\
combined & \texttt{ccicciic\,cccii\,icc} & \textbf{.1808}/\textbf{.0537} & \textbf{.1962}/\textbf{.0701} & \textbf{.2178}/.0857 & \textbf{.2386}/.1033 & \textbf{.2590}/\textbf{.1237} & 1.51 i / 9.25 c \\
\bottomrule
\end{tabular}
\end{table}

\subsection{What the battery actually shows}
\label{sec:act1result}

The result is not the flat null we expected, and the shape it does have is more informative. Three
statements, in decreasing order of how well they are supported.

\paragraph{(a) Every published elicitation method we ran lands below a well-chosen fixed order.}
At sixteen questions on the certified tower, in one contract: our greedy items-only sequence
$.2447/.1017$; entropy$_0$, a 2008 heuristic~\cite{rashid2008learning}, $.2303/.0924$;
CLiME~\cite{liu2025clisk}, a 2025 conversational bandit, $.2218/.0897$; DRE~\cite{kweon2020deep},
the peer-reviewed learned-static bar, $.2185/.0898$; GLM-UCB~\cite{zhao2013interactive}
$.2161/.0833$; Thompson $.2135/.0878$; Golbandi's interview tree~\cite{golbandi2011adaptive}
$.2117/.0844$ (at eight questions; it has no sixteen-question row). The cleanest way to say this is
not about adaptivity at all: \emph{published elicitation methods of every kind --- learned static,
classical adaptive, modern conversational bandit --- fall below a fixed order chosen by greedy
search on the recommender they are scored with.} That the strongest published \emph{learned static}
method is beaten on full accuracy by a 2008 entropy heuristic ($-.0118$, $[-.0153,-.0083]$) is the
sharpest single instance --- though note the same gap is \emph{not} separated on the tail, so the
two axes disagree about that particular pair.

The comparator that carries the claim is therefore the greedy order and not the heuristic. Against
entropy$_0$ specifically the picture is mixed: IGCN separably beats it at one, two and four
questions ($+.0050$, $+.0095$, $+.0113$ full). That is worth stating precisely, because it is
\emph{not} a win for adaptivity --- IGCN's personalised step is active on exactly $0.0\%$ of turns
at those budgets (see (d)), so what beats entropy$_0$ early is IGCN's global information-gain
\emph{static} order.

\paragraph{(b) The acquisition families lose everywhere, and not for want of branching.}
Against the strongest lit-anchored static arm, GLM-UCB, Thompson sampling and the latent-only
information gain lose at \emph{every} budget on \emph{both} axes --- thirty of thirty comparisons,
all with paired intervals excluding zero. The best published conversational bandit loses all thirty
of its comparisons against all three static arms. And the loss is not a failure to adapt: Thompson
gives all $10{,}000$ users a distinct interview by the fourth question and still loses. This is the
part of the result that is robust, contract-invariant (it survives the floored arm unchanged), and
supported by four independent method families.

\paragraph{(c) One method beats the weak static baseline, and still loses to the strong one.}
Golbandi's interview tree~\cite{golbandi2011adaptive} loses to entropy$_0$ at one and two questions
($-.0078$ and $-.0092$, both separated) and beats it at four and eight ($+.0032$ and $+.0064$ full,
$+.0052$ tail at eight). Against the greedy items-only order it loses full NDCG at \emph{every}
budget ($-.0203$, $-.0270$, $-.0174$, $-.0119$, all separated) and loses the tail at one and two
questions, with the tail cells at four and eight failing to separate in either direction
($-.0026$ $[-.0055,+.0003]$ and $-.0016$ $[-.0049,+.0017]$). Against the combined order it loses
both axes at every budget, the tail at eight by $-.0189$ $[-.0226,-.0154]$. Since entropy$_0$ is a
heuristic and the greedy order is the arm actually built to maximise the metric, the summary is the
narrow one: \textbf{no published adaptive selector in this battery beats the strong static order on
either axis at any budget}, and a win over the weak one is not a crossover.

\paragraph{A retraction, in place.} An earlier version of this section reported that the tree took
the tail at eight questions against the strongest static order ($+.0068$, CI $[+.0039,+.0098]$).
That interval was measured against an items-only sequence \emph{selected} under a narrower candidate
mask than the one it was scored under; its tail at eight questions reads $.0776$, where the sequence
rebuilt under the scoring contract reads $.0860$. Re-differenced against the rebuilt order the same
cell is $-.0016$ and does not separate. The tree's own numbers did not move --- the comparator did.
This is worth printing rather than silently correcting for the reason
Section~\ref{sec:positioning} gives: comparator strength is precisely the axis this chapter argues
the literature has left unattended, and the effect is large enough to reverse a separated interval
in our own results.

\paragraph{The claim, and the claim we are careful not to make.} Every published elicitation method
we could replicate lands below a well-chosen fixed order, at every budget, on one contract. What that
does \emph{not} establish is anything about adaptivity. A fixed order is an adaptive policy that
ignores its inputs, so the optimal adaptive policy is at least as good as the optimal fixed one, and
no measurement can put adaptivity below staticness --- only particular methods. The result is a
statement about those methods, and the next two paragraphs are our attempt to say precisely which
property of them it is a statement about.

\paragraph{(d) And a fourth statement, which we did not anticipate and which explains a good deal of
(b): several named adaptive selectors do not condition on the answers at all in this regime.}
Building them was what revealed it, and the mechanism differs in each case. \textbf{IGCN} switches to
its personalised step after five \emph{ratings}; at an answer rate near $0.3$ per question it is
$98\%$ static at eight questions and \emph{exactly} the global-information-gain static order at four
or fewer. \textbf{BALD} is an information-theoretic criterion whose expected-conditional-entropy
term is constant under a homoscedastic linear-Gaussian likelihood, so it reduces algebraically to a
posterior-variance rule --- and since that variance depends only on \emph{which} questions were
answered and never on \emph{what} was answered, it is not answer-adaptive at all. \textbf{A-optimal
design} contains the questions and not the answers by construction. And the conversational-bandit
schedule fires a key-term turn only twice in sixteen, so at short budgets its three variants are
numerically identical.

The diagnostic that makes this legible is not the number of distinct interviews a selector conducts
but the number it conducts \emph{per distinct answer-state}: a method can branch widely and still be
ignoring what it hears. At the second question BALD faced eleven distinct answer-states and asked
two distinct questions. We therefore report, beside every adaptive row, the fraction of turns on
which the method actually adapted.

This changes what the null means. ``Adaptive selection does not pay'' invites the reply that we
implemented these methods badly. The more accurate and more useful statement is that \emph{several
named adaptive selectors are structurally unable to condition on answers in a cold-start interview}
--- through a warm-up threshold denominated in ratings, an algebraic collapse under the standard
likelihood, a value-free criterion, or a schedule that rarely fires --- and that this is a property
of the methods as published, measurable, and different in each case. Counted across the whole
battery, \textbf{twelve of twenty-three rows are value-free}: the rule responds to \emph{which}
questions the user could answer and never to \emph{what} they answered. If adaptivity means
conditioning on answers, then over half the battery is not adaptive in the sense the word is used.

\paragraph{(e) Every method optimises a surrogate for the quantity it is judged on --- except the
fixed order.} Auditing what each rule actually maximises when it picks the next question gives the
plainest available explanation of (a), and it is not a claim about adaptivity at all.
\textbf{Exactly one of the twenty-one distinct selection rules in this battery optimises held-out
ranking quality, and it is the static greedy order.} The nearest prior work on this specific point is
Hu and Yu~\cite{hu2013interview}, an interview tree whose split loss is general and whose paper lists
a \emph{pairwise rank loss} among the options, evaluated with MAP on sampled negatives (ten per
positive) against a top-$k$-popular static comparator. We do not claim to be first to admit a ranking
objective into an interview tree's split, and Section~\ref{sec:huyu} runs their criterion rather than
citing it. The others maximise information gain about a latent variable~\cite{ma2019eddi} or a
cluster label~\cite{rashid2008learning}, minimise posterior
variance~\cite{houlsby2011bald} or an A-optimality trace~\cite{anava2015budget,sepliarskaia2018spq},
reduce squared error of \emph{rating} prediction~\cite{golbandi2011adaptive,zhou2011functional} or of
profile reconstruction~\cite{kweon2020deep,proissl2024greedyslim}, or maximise a UCB or Thompson
index over a predicted rating~\cite{zhao2013interactive,zhang2020conucb,wang2023conlinucb,liu2025clisk,parapar2021dpe}.
The heuristic orders were tuned against rating-prediction error by their own authors --- Rashid
reports MAE~\cite{rashid2002getting,rashid2008learning} and Golbandi
RMSE~\cite{golbandi2011adaptive} --- and transfer to a ranking metric as priors rather than as
certified numbers.

We are careful about what follows from this. Optimising a surrogate is not a defect: expected
NDCG@10 under a neural fold is not tractable inside a bandit index or a tree split, and reaching for
information gain is the reasonable engineering response. The question that decides fairness is not
what a method optimises but what its authors \emph{claimed}, and we audited that separately. Of the
thirteen published methods, \textbf{six put a held-out ranking metric in their primary table and
claimed a ranking improvement} --- DRE ($+11.6\%$ NDCG@10)~\cite{kweon2020deep}, Greedy
SLIM (NDCG@10 on ML-25M, with the same long-tail definition we use)~\cite{proissl2024greedyslim},
RBMF~\cite{liu2011wisdom}, DPE~\cite{parapar2021dpe}, PEBOL~\cite{austin2024bayesian}, and A-optimal
design in its recommender incarnation~\cite{anava2015budget}. For those, holding them to NDCG@10 is holding them to their own headline, and the
comparison needs no apology.

\textbf{The other seven were never validated on a ranking metric by their authors, and each fails to
transfer in a different way.} Golbandi~\cite{golbandi2011adaptive} and the Rashid
lineage~\cite{rashid2002getting,rashid2008learning} --- which supplies entropy$_0$, HELF,
popularity and IGCN --- report RMSE and MAE; Golbandi's paper goes further and disclaims the
application outright (``a decision-tree recommender delivers a quite high RMSE. This makes
decision-trees undesirable as a traditional recommendation system''), which makes it an instrument
paper rather than a recommender paper. The conversational-bandit
trio~\cite{zhang2020conucb,wang2023conlinucb,liu2025clisk} prices \emph{cumulative regret}
over thousands of rounds against re-drawn arm sets: there is no cold user, no interview budget, and no
held-out target anywhere in their protocol, so our rows are ports of a selection rule into a setting
its authors never evaluated. EDDI~\cite{ma2019eddi} is not a recommender method at all --- its
datasets are UCI, MNIST, MIMIC and NHANES --- and enters here as an imported
active-feature-acquisition function. We label
these rows accordingly rather than reporting them as underperforming recommenders.

The precedent here is a reproducibility one~\cite{dacrema2019progress}. A recent benchmark
study~\cite{proissl2024greedyslim} working on this dataset and
this metric --- ML-25M, NDCG@10, the same long-tail split --- set out to include the RMSE-era
interview methods~\cite{golbandi2011adaptive,sepliarskaia2018spq} and could not: ``For Golbandi et
al.\ and Sepliarskaia et al.\ we, unfortunately,
obtained poor results. We therefore decided to not include them in our official experiments.'' Their
own diagnosis follows a sentence later: ``the reproducibility issue within the recommender systems
community is also an issue when it comes to questionnaires. \emph{We had to implement all the
baselines we wanted to consider on our own.}'' We read that as a statement about the availability of
working implementations rather than about the methods themselves, and note that the study is an
unrefereed preprint.

The consequence we draw is correspondingly narrow, and does not require interpreting anyone's
failure: \textbf{the RMSE-era interview methods have not been successfully evaluated on a ranking
metric by anyone, including a group that attempted it on this dataset and this metric.} That is the
gap this battery fills, and it is why each port here is gated on reproducing its published
construction natively \emph{before} it is scored on our ruler --- with the gate's failures reported
as failures rather than quietly scored: EDDI's native snap did not pass, and Greedy SLIM's passed
only in part.

So the contribution is narrower and firmer than ``the field optimises the wrong thing'': the
surrogate-to-ranking transfer was \emph{demonstrated} by six of thirteen, \emph{assumed} by the other
seven, and this battery is the first place that assumption is measured on one contract.

What the audit does establish is that our comparator is not merely \emph{a} strong static order but
the only arm in the battery built by directly maximising the reported metric, and that a reader
should discount the margin accordingly. It also identifies the missing experiment: an adaptive
policy built with the \emph{same} objective and the same optimiser, so that conditioning is the only
thing that varies. Section~\ref{sec:act3} runs it.

\paragraph{What would falsify this, and what currently limits it.} The reading above would be
overturned by an adaptive policy that beats the greedy order on full NDCG at any budget. No
\emph{published} selector in this battery does, and the two nearest misses are instructive rather
than encouraging: IGCN separably beats entropy$_0$ at one, two and four questions while its
personalised step is active on exactly $0.0\%$ of turns, so what wins is its
global-information-gain \emph{static} order; and Golbandi's tree beats entropy$_0$ at depth while
losing to greedy at every budget.

The honest limit is that no arm here isolates the variable of interest. Every comparison confounds
\emph{conditioning on answers} with \emph{choice of objective}, because the only arm optimising the
reported metric is also the only arm that is static. Two experiments separate them, and both are run.
The first partitions users into groups, runs the identical greedy search with the identical objective
inside each, and applies each group's order to held-out users of that group
(Section~\ref{sec:groupgreedy}); static is nested inside that construction, so the whole question is
one of \emph{generalisation}. The second is the recursive form of the same idea --- Golbandi's tree
structure with measured held-out NDCG@10 substituted for the squared-error split
(Section~\ref{sec:act3}) --- and it is the arm that carries this chapter.

\paragraph{Two candidate pools, and why the comparison must be made inside one.}
Table~\ref{tab:greedy} was computed under a narrower candidate mask than Table~\ref{tab:act1}: the
greedy environment removes the user's positive fold-in items and whatever they answered
($62.6$ items per user), while the interview harness removes every rated fold-in item, likes and
dislikes alike ($141.2$ per user), which is this project's protocol-level definition of the pool.
Removing more candidates makes the task easier, so the two tables' columns are not comparable, and
the apparent fact that the entropy$_0$ arm reaches $.2303$ where the greedy combined arm reaches
$.2115$ compares one row from each. Re-scoring both fixed sequences under each mask settles it:
under the wider pool the empty interview alone moves from $.1314$ to $.1586$, and
\textbf{the greedy sequence beats entropy$_0$ at every budget on both axes under \emph{both} masks}
--- by $+.0122$ full / $+.0148$ tail at $q{=}16$ under the narrow pool, and, for the sequence
rebuilt under the protocol pool itself, $+.0144$ $[+.0112,+.0175]$ full and $+.0093$
$[+.0061,+.0124]$ tail for the items-only order and $+.0287$ full / $+.0313$ tail for the combined
one. The re-scoring reproduces each published table to four decimal places once its own mask is
applied, which is what licenses the comparison below.

Placed in the protocol pool alongside the selectors, the strongest static order is therefore
stronger than the strongest heuristic, and the null of Table~\ref{tab:act1} widens rather than
narrows. Restricting attention to the item bank --- the fair comparison, since none of the adaptive
selectors has access to the concept channel --- at $q{=}16$: greedy items-only $.2447/.1017$,
entropy$_0$ $.2303/.0924$, GLM-UCB $.2161/.0833$, Thompson $.2135/.0878$, two-outcome info-gain $.1917/.0749$. The
best adaptive selector trails the strongest items-only static order by $+.0281$
$[+.0247,+.0314]$ full and $+.0178$ $[+.0145,+.0210]$ tail, twice its deficit against the
heuristic. That sequence is now \emph{selected} under the same mask it is scored under, so the
figure is a like-for-like comparison rather than the lower bound earlier drafts reported.

\paragraph{This is not a housekeeping point: a published adaptive-degradation result is
manufactured by it.} The most direct evidence we have that candidate and target formation are
load-bearing comes from outside our own harness. A recent preprint reports, on this very dataset
and metric, that an adaptive questionnaire's accuracy \emph{falls} past roughly ten questions while
a static one keeps rising --- and, to its authors' credit, attributes the fall to their own protocol
rather than to adaptivity. We rebuilt their setup and ran the ablation their text implies. The
decline is indeed an artefact of protocol, but of a different part of it than they name, and the
two parts are separable.

They blame the rule forbidding a questionnaire to recommend the items it asked about, i.e. a
restriction of the \emph{candidate pool}. In their setup that pool is already restricted to items
the user has not rated, so an \emph{answered} item is excluded regardless; their extra rule removes
only asked-but-unanswered items, which offline are guaranteed non-targets. Removing non-targets from
a candidate pool is close to free --- and indeed switching their rule off makes the decline
\emph{steeper} ($-46\%$ against $-14\%$), which cannot happen if that rule is the cause.

What actually shrinks is the \emph{target set}. Their design gives a test user no held-out split, so
the same ratings serve as the answers and as the ground truth: every answered question converts an
item the metric could have credited into an item it can no longer see. Under an arm-independent
held-out split the decline disappears completely and the adaptive arm rises monotonically
($.1637 \to .1761 \to .1857 \to .1936$), while a popularity control declines only $-.012$ against
the adaptive arm's $-.045$ at a matched answer rate --- so the effect is carried by adaptivity, as
they said, and removed by the protocol, as they suspected.

The general form is the rule we adopt: \textbf{the ground truth must not depend on what was asked.}
A protocol that violates it penalises exactly the arms that ask well, and the penalty compounds with
budget, so it is most severe for the methods a study is usually trying to promote. We cite that
paper as motivation rather than as a counterexample: it observed something real, diagnosed it as
protocol rather than blaming the method, and its residual mis-attribution is precisely the
distinction this section exists to draw.

\begin{table}[t]
\centering
\caption{Adaptive selectors versus a static order, consumer = the certified tower, $10{,}000$ test
users, NO-FLOOR, full/tail NDCG@10. MostPop prior $0.1626/0.0262$. The static arm wins every cell
of this table; Section~\ref{sec:flip} reports the consumer on which that reverses. Not on the same
scale as Table~\ref{tab:greedy} --- see the contract note above.}
\label{tab:act1}
\small
\setlength{\tabcolsep}{3pt}
\begin{tabular}{lccccc c}
\toprule
selector & $k{=}1$ & $k{=}2$ & $k{=}4$ & $k{=}8$ & $k{=}16$ & answered@16 \\
\midrule
static (entropy$_0$) & \textbf{.1635/.0433} & \textbf{.1706/.0528} & \textbf{.1831/.0636} & \textbf{.2053/.0792} & \textbf{.2303/.0924} & 5.52 \\
GLM-UCB & .1606/.0400 & .1635/.0452 & .1756/.0578 & .1902/.0673 & .2161/.0833 & 5.24 \\
Thompson & .1627/.0376 & .1676/.0437 & .1748/.0551 & .1908/.0677 & .2135/.0878 & 3.39 \\
latent-only info-gain$^\dagger$ & .1537/.0336 & .1539/.0373 & .1618/.0432 & .1747/.0611 & .1917/.0749 & 3.84 \\
\bottomrule
\end{tabular}
\end{table}

The ordering repeats when the consumer is changed to RecVAE, to TaNP, or to EDDI's own Partial-VAE
--- including the self-consistent cell in which EDDI selects and EDDI consumes ($.1493/.0488$
against $.1769/.0480$ for the static order at sixteen questions), so on these four consumers the
result is not an artefact of pairing a selector with a recommender it was not designed for.
Section~\ref{sec:flip} reports the fifth consumer, on which it reverses.
All thirty comparisons in Table~\ref{tab:act1} separate under a paired bootstrap ($10{,}000$
resamples over the per-user difference, the same $10{,}000$ users in both arms, self-comparison
returning exactly zero): full-axis deltas against the static order run from $-.0029$ to $-.0137$
for GLM-UCB, $-.0022$ to $-.0177$ for Thompson, and $-.0098$ to $-.0386$ for the latent-only
info-gain arm as the budget grows from one question to sixteen.

\subsection{The battery so far, in one contract}
\label{sec:battery}

\paragraph{One scope statement governs every table below: this literature asks about items.}
Golbandi's tree, DRE, RBMF, RMVA, IGCN, fMF, EDDI, the interactive-collaborative-filtering bandits
and the recent SLIM-based questionnaires all select \emph{items} to ask about, and so do all six of
our literature-anchored static arms and all three of our adaptive selectors. The exceptions are
narrow and we name them because they are the only available rivals for the coarse channel: the
conversational-bandit family, which asks about a key-term or an item in one algorithm, and the
topic bandit. Consequently \textbf{the adaptive-versus-static comparison of Act 1 is a result about
the item channel}, and the correct comparator for those selectors is the greedy \emph{items-only}
sequence, not the combined one, which draws on a channel no adaptive selector here can reach. Where
the combined sequence appears it is a bound on what selection can buy given both channels, never a
like-for-like opponent. The concept channel's own comparison is a separate claim with its own
comparators, and conflating the two would repeat, at the level of the action space, exactly the
contract error the candidate-pool discussion above warns about.

Table~\ref{tab:battery} collects every selection method measured to date under a single information
contract, and Table~\ref{tab:coverage} states what is still missing. Both are working-draft
inventories: the point of printing them is that the empty cells are as visible as the full ones.

\begin{table}[t]
\centering
\caption{Every selector measured to date. One contract throughout: consumer = certified tower,
protocol candidate pool (every rated fold-in item removed), NO-FLOOR, behavioural answers, all
$10{,}000$ test users, no sampling. full/tail NDCG@10. Answer-free reference (popularity prior):
$.1626/.0262$. The three greedy rows are the sequences rebuilt under this table's own candidate mask
(Section~\ref{sec:act3}); selection contract and scoring contract now agree, so they are no longer
lower bounds. $^\S$ marks a \emph{value-free} rule: it responds to which questions turned out to be
answerable and never to what was answered (Section~\ref{sec:act1result}, statement (d)).
$^\ddagger$ The same reduced acquisition as the row above it, under a three-outcome response model
(no answer / like / dislike) rather than a two-outcome one; it is no more the published criterion
than that row is, and is carried for the mechanism comparison rather than as a replication.}
\label{tab:coverage-note}
\label{tab:battery}
\footnotesize
\setlength{\tabcolsep}{3pt}
\begin{tabular}{llccccc c}
\toprule
selector & type & $k{=}1$ & $k{=}2$ & $k{=}4$ & $k{=}8$ & $k{=}16$ & ans@16 \\
\midrule
\multicolumn{8}{l}{\emph{optimised static (this chapter's comparator)}}\\
greedy combined (item+concept) & static & \textbf{.1808/.0537} & \textbf{.1962/.0701} & \textbf{.2178}/.0857 & \textbf{.2386}/.1033 & \textbf{.2590/.1237} & 1.5i/9.2c \\
greedy items-only & static & .1760/.0434 & .1884/.0563 & .2038/.0677 & .2235/.0860 & .2447/.1017 & 3.11 \\
greedy concepts-only & static & \textbf{.1808/.0537} & \textbf{.1962/.0701} & .2141/\textbf{.0858} & .2347/\textbf{.1074} & .2447/\textbf{.1161} & 12.36 \\
\midrule
\multicolumn{8}{l}{\emph{literature-anchored static heuristics}}\\
entropy$_0$~\cite{rashid2008learning} & static & .1635/.0433 & .1706/.0528 & .1831/.0636 & .2053/.0792 & .2303/.0924 & 5.52 \\
popularity~\cite{rashid2002getting} & static & .1606/.0400 & .1647/.0470 & .1746/.0579 & .1943/.0686 & .2228/.0848 & 5.42 \\
HELF~\cite{rashid2008learning} & static & .1603/.0354 & .1628/.0420 & .1755/.0588 & .1892/.0725 & .2157/.0838 & 3.34 \\
pure entropy~\cite{rashid2002getting} & static & .1582/.0319 & .1578/.0327 & .1580/.0323 & .1586/.0317 & .1599/.0357 & 0.08 \\
random (bank) & static & .1586/.0310 & .1588/.0311 & .1588/.0316 & .1594/.0329 & .1607/.0353 & 0.12 \\
random (catalogue) & static & .1585/.0310 & .1588/.0310 & .1589/.0318 & .1593/.0315 & .1612/.0341 & 0.12 \\
\midrule
\multicolumn{8}{l}{\emph{adaptive}}\\
GLM-UCB~\cite{zhao2013interactive} & adaptive & .1606/.0400 & .1635/.0452 & .1756/.0578 & .1902/.0673 & .2161/.0833 & 5.24 \\
Thompson~\cite{zhao2013interactive} & adaptive & .1627/.0376 & .1676/.0437 & .1748/.0551 & .1908/.0677 & .2135/.0878 & 3.39 \\
latent-only info-gain$^\dagger$~\cite{ma2019eddi} & adaptive & .1537/.0336 & .1539/.0373 & .1618/.0432 & .1747/.0611 & .1917/.0749 & 3.84 \\
EDDI three-outcome$^\ddagger$~\cite{ma2019eddi} & adaptive & .1607/.0363 & .1597/.0388 & .1732/.0497 & .1869/.0660 & .2041/.0834 & 2.86 \\
\midrule
\multicolumn{8}{l}{\emph{published methods, snap-validated and ported}}\\
Golbandi tree (WSDM'11)~\cite{golbandi2011adaptive} & adaptive & .1557/.0351 & .1614/.0478 & .1864/.0650 & .2117/.0844 & --- & 2.75 \\
IGCN (2008)~\cite{rashid2008learning} & adaptive & .1685/.0407 & .1800/.0569 & .1944/.0701 & .2031/.0747 & .2239/.0912 & 4.19 \\
DRE (WWW'20)~\cite{kweon2020deep} & learned static & .1658/.0396 & .1693/.0465 & .1763/.0586 & .2074/.0789 & .2185/.0898 & 4.6 \\
CLiME (KDD'25)~\cite{liu2025clisk} & adaptive & .1606/.0400 & .1634/.0452 & .1776/.0549 & .1914/.0701 & .2218/.0897 & 4.80 \\
ConLinUCB-BS-static~\cite{wang2023conlinucb} & static (concepts) & .1644/.0374 & .1659/.0402 & .1724/.0468 & .1829/.0598 & .1891/.0700 & 0.00 \\
ConUCB (WWW'20)~\cite{zhang2020conucb} & adaptive & .1584/.0346 & .1591/.0364 & .1651/.0447 & .1805/.0574 & .1874/.0699 & 1.80 \\
DPE-Hellinger (WSDM'21)~\cite{parapar2021dpe} & adaptive & .1627/.0398 & .1680/.0480 & .1817/.0610 & .2002/.0765 & .2192/.0925 & 3.9 \\
DPE-ILD (WSDM'21)~\cite{parapar2021dpe} & adaptive & .1627/.0398 & .1621/.0401 & .1613/.0397 & .1609/.0390 & .1603/.0375 & 0.32 \\
BALD (2011)$^\S$~\cite{houlsby2011bald} & value-free & .1754/.0448 & .1853/.0473 & .1934/.0572 & .2057/.0683 & .2203/.0824 & 2.80 \\
A-optimal (WWW'15)$^\S$~\cite{anava2015budget,sepliarskaia2018spq} & value-free & .1693/.0498 & .1724/.0523 & .1702/.0549 & .1694/.0544 & .1740/.0563 & 1.81 \\
\midrule
\multicolumn{8}{l}{\emph{ceiling --- reads the user's own ratings at selection time; NEVER a baseline}}\\
clairvoyant oracle & --- & .1893/.0680 & .2194/.0965 & .2544/.1281 & .2930/.1597 & \textbf{.3292/.1927} & 12.9 \\
\bottomrule
\end{tabular}
\end{table}

Read down the table, with paired bootstrap intervals attached (10{,}000 resamples over the per-user
difference; a selector compared with itself returns exactly zero). The optimised static comparator
leads every adaptive selector at every budget, and every one of those thirty comparisons separates
--- and that holds against all three rebuilt static orders, thirty of thirty each. Widened to every
adaptive arm measured under this contract --- the table's thirteen adaptive rows plus the
three-outcome sub-variant discussed in Section~\ref{sec:answerability}, $138$ cells in all --- the
combined and concepts-only orders lead and separate in all $138$; the items-only order leads in
$136$ and separates in $134$, and the four exceptions are all tail cells that fail to separate in
either direction (Golbandi's tree at four and eight questions, IGCN at two and four). Neither
three-outcome variant adds an exception.
The three degenerate arms are degenerate for a diagnosable reason --- they collect roughly one tenth
of an answer per interview --- and they are mutually indistinguishable, as they should be. The
middle of the table is less ordered than it looks: the two bandit selectors are mostly
indistinguishable from HELF, while the latent-only info-gain arm sits significantly \emph{below}
all three conventional heuristics ($-.0240$ against HELF and $-.0311$ against popularity at
$k{=}16$, both separating). Among the static arms, entropy$_0$ separates above both popularity and
HELF everywhere, but popularity and HELF do not separate from each other at one, two or four
questions and HELF wins the tail at eight --- so the conventional ladder is entropy$_0$ above an
unordered pair, not a total order. The gap that matters for Act 1 is the item-bank one at $k{=}16$:
greedy items-only $.2447$ against the best adaptive selector's $.2161$.

\begin{table}[t]
\centering
\caption{Battery coverage. \cmark{} = measured in the contract of Table~\ref{tab:battery};
\xmark{} = not yet. ``Native'' = run with the recommender the method was published with; ``own
static'' = compared against the static baseline its own paper used. $^\P$ DPE's two published
acquisition variants are the rows built; that paper reports no scalar numbers and no static
comparator of its own, so ``native'' and ``own static'' are not defined for it.}
\label{tab:coverage}
\footnotesize
\setlength{\tabcolsep}{4pt}
\begin{tabular}{lllcccc}
\toprule
method & source & type & built & ours & native & own static \\
\midrule
greedy prefix-optimal & this chapter & static & \cmark & \cmark & n/a & n/a \\
NDCG tree (Section~\ref{sec:act3}) & this chapter & adaptive & \cmark & \cmark & n/a & n/a \\
entropy$_0$ / HELF / pop / random & \cite{rashid2002getting,rashid2008learning} & static & \cmark & \cmark & \xmark & n/a \\
GLM-UCB, Thompson & \cite{zhao2013interactive} & adaptive & \cmark & \cmark & \xmark & \xmark \\
latent-only info-gain$^\dagger$ & \cite{ma2019eddi} & adaptive & \cmark & \cmark & \cmark & \xmark \\
EDDI three-outcome$^\ddagger$ & \cite{ma2019eddi} & adaptive & \cmark & \cmark & \xmark & \xmark \\
\midrule
ConLinUCB-MCR / -BS / ConUCB & \cite{wang2023conlinucb,zhang2020conucb} & adaptive & \cmark & \cmark & \xmark & \xmark \\
CLiME & \cite{liu2025clisk} & adaptive & \cmark & \cmark & \xmark & \xmark \\
ConLinUCB-BS-static & \cite{wang2023conlinucb} & static (concepts) & \cmark & \cmark & \xmark & \xmark \\
Golbandi tree (split search) & \cite{golbandi2011adaptive} & adaptive & \cmark & \cmark & \cmark & \cmark \\
Hu \& Yu split criteria & \cite{hu2013interview} & adaptive & \cmark & \cmark & \cmark & \xmark \\
ICF (4 policies) & \cite{zhao2013interactive} & adaptive & \xmark & \xmark & \xmark & \xmark \\
IGCN & \cite{rashid2008learning} & adaptive & \cmark & \cmark & \cmark & \xmark \\
DPE-Hellinger / DPE-ILD$^\P$ & \cite{parapar2021dpe} & adaptive & \cmark & \cmark & n/a & n/a \\
BALD & \cite{houlsby2011bald} & value-free & \cmark & \cmark & \xmark & \xmark \\
SEL-BALD & \cite{gao2024selbald} & adaptive & \xmark & \xmark & \xmark & \xmark \\
A-optimal design & \cite{anava2015budget,sepliarskaia2018spq} & value-free & \cmark & \cmark & \xmark & \xmark \\
PEBOL acquisition & \cite{austin2024bayesian} & adaptive & \xmark & \xmark & \xmark & \xmark \\
fMF & \cite{zhou2011functional} & adaptive & \xmark & \xmark & \xmark & \xmark \\
PERE & \cite{nguyen2024pere} & adaptive & \xmark & \xmark & \xmark & \xmark \\
DRE & \cite{kweon2020deep} & learned static & \cmark & \cmark & \cmark & \xmark \\
Greedy SLIM / QBandit & \cite{proissl2024greedyslim,christakopoulou2016towards} & static/adaptive & \xmark & \xmark & \xmark & \xmark \\
RBMF (maxvol) / RMVA & \cite{liu2011wisdom,fonarev2016rectangular} & static & \xmark & \xmark & \xmark & \xmark \\
GreedyExtend & \cite{golbandi2010bootstrapping} & static & \xmark & \xmark & \xmark & \xmark \\
The FacT / FacT-CRS & \cite{tao2019fact,haque2022factcrs} & adaptive & \xmark & \xmark & \xmark & \xmark \\
UpsRec (privileged oracle) & \cite{lin2023upsrec} & adaptive & \xmark & \xmark & \xmark & \xmark \\
MeLU candidate score & \cite{lee2019melu} & static bank & \xmark & \xmark & \xmark & \xmark \\
a-stratification & \cite{chang1999astratification} & static bank & \xmark & \xmark & \xmark & \xmark \\
\bottomrule
\end{tabular}
\end{table}

\paragraph{The published methods, measured.} Two families have since been snap-validated against
their own papers and run here. \textbf{Golbandi's interview tree}~\cite{golbandi2011adaptive} is the
only published selector in this battery to beat \emph{a} static comparator: on our tower it loses to
entropy$_0$ at one and two questions ($-.0078$, $-.0092$, both separated) and wins at eight
($+.0064$ full, $+.0052$ tail, separated). Against the rebuilt greedy orders it wins nothing at any
budget on either axis. \textbf{The conversational-bandit
family}~\cite{zhang2020conucb,wang2023conlinucb,liu2025clisk} ---
the only published line that asks about a concept or an item in one algorithm, with no target item
and no simulator --- tops out at $.2218/.0897$ (CLiME at sixteen questions), above our own bandit
reimplementations but below both entropy$_0$ ($-.0085$) and the strongest static ($-.0231$).

Two features of that family are worth stating because they are properties of the published
method, not of our port. Its own key-term schedule fires on turns three and eight only, so
\textbf{sixteen questions spend two of them on concepts} and the one- and two-question rows contain
no concept questions at all --- the method is item-dominated by construction. And its pure concept
arm, whose order is the barycentric spanner of the key-term features, is \textbf{the first
published opponent our concept channel has had}. Our bank beats it by a wide, separated margin at
every budget on both axes; the size depends on which of two answer conventions both arms are scored
under, and we report both rather than pick: $+.0407$ $[+.0381,+.0432]$ full and $+.0276$
$[+.0246,+.0308]$ tail at sixteen questions under the convention the bandit harness uses, and
$+.0553$ $[+.0524,+.0582]$ full and $+.0408$ $[+.0378,+.0440]$ tail under the convention the
concept channel was trained on. \REDO{the conversational-bandit family produced all of its concept
answers under the convention it was \emph{not} trained on, which costs a concept-only order up to
$.0177$ tail --- so its concept numbers are floors, and the family should be re-run before any of
them is quoted as an estimate}

\PLAN{the remaining \xmark{} rows of Table~\ref{tab:coverage}, ordered by leverage: the second
learned static bar (Greedy SLIM, the only published protocol sibling of DRE), then the remaining
information-theoretic arms (SEL-BALD, PEBOL's acquisition), then fMF, PERE and RBMF. Also owed: the
two remaining key-term representations, which together adjudicate how a concept should be
represented.}

\subsection{The loss is not a failure to branch}

The obvious objection --- ``your adaptive policies simply did not adapt'' --- is refuted by the
harness's own diagnostic. Counting distinct question-prefixes over the $10{,}000$ users:

\begin{table}[h]
\centering
\caption{Distinct interviews actually conducted (ask-prefixes over $10{,}000$ users).}
\label{tab:prefixes}
\small
\begin{tabular}{lccccc}
\toprule
selector & $k{=}1$ & $k{=}2$ & $k{=}4$ & $k{=}8$ & $k{=}16$ \\
\midrule
static (entropy$_0$) & 1 & 1 & 1 & 1 & 1 \\
GLM-UCB & 1 & 3 & 23 & 335 & 5{,}739 \\
Thompson & 759 & 8{,}083 & 10{,}000 & 10{,}000 & 10{,}000 \\
latent-only info-gain$^\dagger$ & 1 & 2 & 7 & 99 & 1{,}993 \\
EDDI three-outcome$^\ddagger$ & 1 & 2 & 8 & 112 & 1{,}217 \\
\bottomrule
\end{tabular}
\end{table}

Thompson sampling conducts a distinct interview for every user from the fourth question onward and
still loses to one order asked of everybody. This is \emph{divergence without exploitability},
measured directly rather than inferred, and it is the honest form of the mechanism claim in
Section~\ref{sec:act2}: manufacturing per-user variety is not the same as having something to
adapt to.

\subsection{What the adaptive selectors spend their budget on}
\label{sec:answerability}

The answered column of Table~\ref{tab:act1} carries part of the mechanism. At sixteen questions the
static order collects $5.52$ answers; Thompson collects $3.39$ and the two-outcome info-gain arm
$3.84$ (its three-outcome variant, $2.86$). Both criteria
prefer questions with high posterior movement or high score uncertainty, and on an implicit-feedback
catalogue those are precisely the items almost nobody has seen, so they spend budget on questions
the user cannot answer. This is the same asymmetry that separates the two question channels on this
stack --- sixteen item questions yield $3.11$ answers where sixteen concept questions yield $12.36$
(Table~\ref{tab:greedy}) --- reappearing as a property of the \emph{policy} rather than of the bank.
Chapter A reports the channel version of this figure at a budget of eight; the numbers quoted here
are this chapter's own, measured on the certified pair.

We state the criterion-side conclusion carefully, because the obvious strong form is not ours to
make. Answerability-weighted acquisition is not new: expected information gain multiplied by an
estimated probability that the user can answer appears in Rashid's popularity-times-entropy
~\cite{rashid2002getting}, in a granted patent, in an ``asking easy questions'' formulation that
subtracts the human's own answer entropy~\cite{biyik2023soft}, in region-elicitation weighted by an
estimated experience probability~\cite{nguyen2024pere}, and in recent active learning with instance
rejection~\cite{gao2024selbald}. What we add is not the
weighting but the observation that it is a reason to change \emph{channel}: prior work prices
answerability within one fixed question space, whereas the coarse and fine channels here differ in
answer rate by a factor of four, and the loss-optimal channel changes with placement
(Section~\ref{sec:act4}). The first of the two info-gain rows uses a two-outcome response model,
under which ``has not seen it'' and ``dislikes it'' are the same observation, so its acquisition
cannot represent --- and therefore cannot price --- the event that no answer arrives; the second row
gives it that third outcome, and the paragraph below reports what changes.

It is not the whole mechanism, and the same table says so: GLM-UCB collects $5.24$ answers,
essentially the static order's rate, and still loses at every budget. Unanswerability explains the
two selectors that starve themselves; it does not explain the one that is fed and still does worse.
The residual explanation must be that the questions it chooses, while answerable, are worth less
than the ones a loss-optimised order chooses --- an optimisation-target mismatch (score uncertainty
against ranking loss) rather than an answerability failure. \REDO{separate the two effects
directly: an answerability-matched control in which each adaptive selector's bank is restricted to
questions of the static arm's answer rate, so the remaining gap is attributable to the criterion
alone}

\paragraph{$^\dagger$ A validation failure we report rather than absorb.}
The information-gain row above is \emph{not} the published EDDI acquisition, and we relabel it
accordingly. Validating our implementation against the original authors' own tabular benchmark
showed that the reduced criterion we had implemented --- the latent-only term, which measures how
much an answer moves the posterior over the latent state --- ranks \emph{below random ordering} on
their data, while their full criterion, which additionally asks how much that movement changes the
prediction one actually cares about, reproduces their published ordering. The reduction was made on
a stated but incorrect worry about leaking evaluation targets; the discarded term samples its
hypothetical answers from the model and observes nothing held out. Until the full criterion is
implemented, this row measures our reduction and not their method, and no conclusion about
information-theoretic acquisition should be drawn from it. We keep it in the table because a
reported validation failure is more informative than a silently omitted arm.

A second, more interesting consequence follows. We had expected that giving the criterion a third
possible response --- \emph{no answer arrives}, distinct from a dislike --- would let it price
answerability on its own, since a question the user cannot answer carries no information by
construction. It does. Measured on the same contract, the three-outcome variant collects $2.86$
answers in sixteen questions against the two-outcome arm's $3.84$ --- a quarter fewer --- and ranks
better from them: the paired in-run difference is positive and excludes zero at every budget on full
accuracy ($+.0070$, $+.0057$, $+.0114$, $+.0122$, $+.0124$ at one, two, four, eight and sixteen
questions) and at four of five budgets on the tail, the two-question cell being a null
($+.0015$, $[-.0001,+.0031]$). Pricing the event that no answer arrives is worth roughly a quarter
of this arm's deficit against the loss-optimised items-only order on full accuracy ($-.0530$ to
$-.0406$ at sixteen questions) and about a third of it on the tail ($-.0268$ to $-.0183$), and it
closes neither. The likely mechanism is structural, and we state it as the code's contract rather
than as a measurement: the two-outcome arm weights its information term by a softmax across items,
which is a relative next-draw probability and not the per-variable marginal that a factorised
likelihood asks for, so the weight behaves like a popularity ranking --- and popularity is a good
proxy for answerability. The three-outcome arm replaces that weight with two proper per-item
quantities in $[0,1]$, which lets the information term drive selection --- and an information term
with no target to serve prefers exactly the questions that move the posterior most, which are the
ones users cannot answer. That is what the answer count records, and it is why the repaired arm is
still a long way behind a loss-optimised order. \textbf{Answerability is not
something an uninformed information criterion recovers for free; it has to be either priced
explicitly or inherited from a target-coupled objective.}

\paragraph{$^\ddagger$ The same validation failure governs the three-outcome row.}
That row repairs the response model over the \emph{same} reduced acquisition. The target term of the
published criterion is implemented nowhere in this chapter; the native snap the reduction fails is
unchanged by the repair (pooled average rank $3.402$ against random ordering's $3.236$ on the
original authors' own benchmark); and the statement above about the two-outcome weight is read off
the implementation, not measured --- the diagnostic that would quantify it has no reproducible
artifact and is not quoted here as a number. Both rows are the partial criterion, and neither is a
replication of the published method: what the pair measures is a response model, holding the
acquisition fixed.

The three-outcome arm admits a second sub-variant, which additionally \emph{conditions} on a dislike
as a signed observation. We do not carry it in the table, for a reason that is not its numbers: the
frozen encoder was trained on likes and has no learned token for a disliked observation, so
conditioning on one is a documented approximation, whereas letting the third outcome enter through
the answerability weight alone is what an absence-is-not-a-value model licenses. Its numbers are
close and do not force the choice --- $.1607/.0363$, $.1607/.0383$, $.1655/.0442$, $.1852/.0649$,
$.2041/.0815$, and $3.16$ answers --- level with the carried row at one and sixteen questions on
full accuracy to within $.0001$, ahead of it only at two questions ($.1607$ against $.1597$), and
behind it on full accuracy at four and eight and on the tail at every budget from two onwards. It is
also three times as expensive to run, branching into $3{,}662$ distinct interviews at sixteen
questions against $1{,}217$, and asking more questions for it. It introduces no exception to the
counts reported for Table~\ref{tab:battery}.

\subsection{The result that complicates the story, reported rather than buried}
\label{sec:flip}

On one consumer the ordering reverses, though only on one axis and only for some selectors, and the
qualification is part of the finding. With Golbandi's leaf model as the consumer --- which matches a
cold user to the train users sharing their like/dislike/unknown pattern --- Thompson sampling beats
the static order on full accuracy at every budget it can be run at, with paired intervals excluding
zero throughout ($+.0100$, $+.0171$, $+.0148$, $+.0281$ at one, two, four and eight questions).
On the tail the reversal is not general: Thompson separates \emph{against} the static order at four
questions ($-.0064$, CI $[-.0095,-.0033]$) while separating for it at one, two and eight. The
latent-only info-gain arm shows the same split at eight questions ($+.0131$ on full,
$-.0034$ on tail, both separating). GLM-UCB does not participate in the reversal at all: $+.0007$
full and $+.0008$ tail, neither separating. So this is a property of the \emph{sampling} selectors
on the \emph{full} axis, not of adaptive selection in general.

The mechanism is not mysterious, and it sharpens rather than weakens the chapter. That consumer is
paid in \emph{patterns}: its estimate improves when users are spread across distinct answer
signatures, so a per-user randomised selector helps it by construction, while a single global order
concentrates most users in a few cells. It is also the one consumer whose input contract includes
the unknowns, so a question the user cannot answer is still informative to it. Which selector wins
therefore depends on which recommender one assumes --- the selection-side appearance of Chapter A's
instrument-dependence result, and a caution that applies to every published comparison in this
literature, including ours.

\subsection{The one arm that clears the bar: the same greedy search, run inside groups}
\label{sec:groupgreedy}

Every comparison above confounds two things. The only arm optimising the reported metric is also the
only arm that is static, so ``conditioning on answers'' and ``choice of objective'' vary together and
the null cannot be attributed to either. This subsection removes the confound in the only way that is
clean: hold the objective, the optimiser, the bank, the scoring contract and the val/test discipline
fixed, and vary \emph{only} what the search is conditioned on. Users are partitioned by favourite
genre, computed from each user's own \emph{fold-in} ratings alone --- the most frequent genre among
their liked fold-in items, with targets excluded (asserted zero leaks out of $10{,}000$ users on both
splits); all $10{,}000$ are partitioned, none dropped and none merged, and the distribution is
dominated by one tag (Drama $5{,}481$, Action $1{,}582$, Comedy $1{,}475$, then a long tail down to
groups of one). The identical greedy search, on the identical objective (cohort-mean held-out full
NDCG@10) with the identical optimiser, is then run separately inside each group; each group's order
is fifteen items and the comparator is the global fifteen-item order rebuilt under the same
protocol pool, compared item-for-item at equal $K$. Static is nested in the construction --- if one
order for everyone were best, each group's search could recover it, and all eighteen groups pass that
containment check on validation. \textbf{The genre answer is treated as free side information}: it
costs no turn and is not folded into the recommender, it only changes which fifteen items are asked,
so every number here is an upper bound. Charging it is not a rounding error --- the sixteenth pick of
the rebuilt static order is worth $+0.0029$ on validation, the same order as the effect itself.

Selection ran on the $10{,}000$ validation users and the disjoint $10{,}000$ test users were looked at
once, under a pre-specified paired percentile bootstrap ($10{,}000$ resamples, seed $98765$), for two
arms fixed before the look: \textsc{all-groups}, which routes every user to their own group's order,
and \textsc{fallback-200}, which routes the $453$ test users whose group had fewer than $200$
validation users back to the global order. At $K{=}15$, \textsc{fallback-200} gains $+0.00862$ full
($[+0.00595,+0.01131]$) and $+0.00591$ tail ($[+0.00326,+0.00862]$); \textsc{all-groups} gains
$+0.01015$ full and $+0.00719$ tail. Against a rebuilt global baseline of $0.24295$ that is
$\sim\!3.5\%$ relative, and \textbf{it is the first arm in this chapter to beat the strong static
order on held-out users}. The profile is not flat: $K{=}1$ is identical by construction (both orders
open with the same item), $K{=}2$ and $K{=}4$ separate on both axes, and the tail stops separating at
$K{=}8$ ($+0.00203$, $[-0.00073,+0.00482]$) before separating again at $K{=}15$. Three qualifications
are the size of the result rather than footnotes to it. \textbf{Drama --- $54.8\%$ of the cohort ---
is a null} ($+0.00150$, $[-0.00187,+0.00486]$): ``favourite genre = Drama'' is the largest and vaguest
tag, a default bucket collecting everyone whose taste is broad, so there is nothing to condition on
inside it and it drags the cohort-weighted number down. \textbf{Validation overstates the effect,
uniformly}: all seventeen non-empty groups shrank from validation to test, not one held its gain, by
$1.80\times$ on the full axis (the tail barely moves, $1.02\times$). And the orders are genuinely
different objects rather than perturbations --- each group's fifteen picks share between zero and six
items with the global fifteen (Table~\ref{tab:firstthree}) --- which makes this a second instance of
divergence \emph{with} exploitability only where the partition is real: gains run from $+0.017$ to
$+0.090$ for Sci-Fi, Adventure, Thriller and Action, while the small groups fitted on tens of users
behave exactly as a free parameter per group should (Horror, fitted on $51$, puts \emph{Shrek 2}
third in a horror interview; Fantasy, fitted on $11$, lands at $-0.116$ on test), which is why
\textsc{fallback-200} exists and was fixed before the test look.

\begin{table}[t]
\centering
\caption{The first three questions of each order. The global order is the comparator; each group's
order is produced by the same search on the same objective, differing only in which users it is
optimising for.}
\label{tab:firstthree}
\footnotesize
\setlength{\tabcolsep}{4pt}
\begin{tabular}{p{1.6cm}p{3.4cm}p{3.4cm}p{3.6cm}}
\toprule
order & question 1 & question 2 & question 3 \\
\midrule
global & Lord of the Rings: The Fellowship of the Ring & Batman (1989) & Star Wars: Episode V --- The Empire Strikes Back \\
\addlinespace
Sci-Fi & The Terminator (1984) & Star Wars: Episode VI --- Return of the Jedi & Alien (1979) \\
Adventure & Lord of the Rings: The Fellowship of the Ring & Harry Potter and the Goblet of Fire (2005) & Harry Potter and the Deathly Hallows: Part 1 \\
Action & Lord of the Rings: The Return of the King & The Terminator (1984) & Star Wars: Episode V --- The Empire Strikes Back \\
Thriller & True Lies (1994) & The Fugitive (1993) & The Rock (1996) \\
\addlinespace
Drama & Clear and Present Danger (1994) & Lord of the Rings: The Fellowship of the Ring & 2001: A Space Odyssey (1968) \\
Horror & From Dusk Till Dawn (1996) & Alien (1979) & \textbf{Shrek 2 (2004)} \\
\bottomrule
\end{tabular}
\end{table}

\paragraph{What this is not.} It is not evidence that adaptive selection works. Every group still asks
a \emph{fixed} order; nothing conditions on anything during the interview. What pays is knowing one
crude fact \emph{before} the interview starts, and it pays only where the fact names a real taste ---
counting how many of a group's fifteen picks carry its own genre tag, against the global order's
congruence on the same genre as a matched control, every group with a lift of $+5$ or more has a
positive test gain and every group with a lift of $+3$ or less, Drama included, is null or negative.
The leverage is at the opener, not in the loop. Section~\ref{sec:act3} puts the conditioning inside
the loop, charges every question, and asks whether it still pays.
\PLAN{charge the genre turn
--- each group's order at $K{-}1$ items against the global order at $K$ --- and, separately, fold the
genre answer through the concept channel instead of discarding it, so that the turn buys a
recommender update as well as a routing decision}


\section{Act 2: the diagnosis --- an objective mismatch}
\label{sec:act2}

Act 1 audited what each rule maximises and found one of twenty-one optimising the reported quantity.
That is a description. This section measures it: how far is each published criterion from the metric
the field reports, at the states an adaptive interview actually visits? And then it asks the
complementary question --- whether there is anything in the users to adapt to at all.

\subsection{How far is the surrogate from the metric? The exact answer, at real states}
\label{sec:surrogategap}

The NDCG tree of Section~\ref{sec:act3} persists, at every node it visits, the exact measured
cohort-mean NDCG@10 marginal gain of \emph{every} candidate in the pool over that node's own users.
That is $22$ nodes $\times$ $\sim\!1{,}000$ candidates of ground truth for the quantity each published
criterion approximates --- already banked, at zero additional cost. We score three published criteria
against it: Golbandi's within-node rating variance~\cite{golbandi2011adaptive}, Hu and Yu's pairwise
rank loss~\cite{hu2013interview}, and the pointwise logistic loss their ranking experiments actually
ran.

\begin{table}[h]
\centering
\caption{Published split criteria against exact measured NDCG@10 gains, at all $22$ nodes of the
grown tree. ``Regret'' is the measured NDCG@10 forgone by following that criterion instead of the
metric, user-weighted; the last row is the expected regret of a uniform random pick, which is the
control a criterion must beat.}
\label{tab:criterionregret}
\small
\begin{tabular}{lrrrrr}
\toprule
criterion & mean $\rho$ & nodes $\rho{>}0$ & argmax hit & top-10 recall & regret \\
\midrule
variance~\cite{golbandi2011adaptive} & $+0.010$ & 13/22 & 1/22 & 0.03 & $+0.01664$ \\
pairwise rank loss~\cite{hu2013interview} & $-0.060$ & 10/22 & 0/22 & 0.01 & $+0.01738$ \\
logistic loss~\cite{hu2013interview} & $-0.104$ & 9/22 & 1/22 & 0.03 & $+0.01282$ \\
\midrule
\emph{control: uniform random pick} & --- & --- & --- & 0.10 & $+0.01955$ \\
\bottomrule
\end{tabular}
\end{table}

Read the last column first. \textbf{Following any of these three criteria forgoes between $66\%$ and
$89\%$ of the NDCG@10 that a uniform random pick would forgo.} The rank correlations are
approximately zero and two of the three are negative, so these are not weak-but-aligned surrogates:
at the states this tree actually visits they are close to uninformative about the quantity the field
reports. Golbandi's criterion picks the measured argmax at $1$ node of $22$; the pairwise loss at
none, with a top-ten that overlaps the metric's top ten \emph{less} often than a random draw.

The root is a real and instructive exception. There the logistic criterion's argmax \emph{is} the
metric's argmax --- the same concept our tree and the static order both open with --- with regret
identically zero, while the variance criterion ranks that question $250$th of $1{,}000$ and the
pairwise loss $75$th. One level down, the logistic criterion's regret is no better than the others'.
So the fair statement is not that these criteria are bad: it is that \textbf{they can agree with the
metric where the cohort is largest and the signal strongest, and stop agreeing exactly where an
adaptive method has to earn its keep.}

\subsection{Is there anything to adapt to? Placement and dispersion}

If adaptive selection is to pay, users must occupy states in which \emph{different} questions are
optimal. We measure that directly, with a loss-coupled criterion: for state $\zz$ and question $q$,
$U(q\mid \zz)=\sum_a P(a\mid \zz)\,V(\zz\oplus(q,a))$, where $V$ is the discounted top-ten mass of
the recommender's own predictive distribution --- i.e.\ the expected NDCG@10 of the slate it would
show. Because the criterion marginalises over the answer, it is not a variance-only criterion whose
dispersion is zero by construction. We then ask how many \emph{distinct} questions are optimal
across a population, reported as the perplexity of the argmax distribution (an effective number of
questions).

\begin{table}[h]
\centering
\caption{Effective number of distinct loss-optimal next questions (perplexity of the argmax
distribution) when the question bank is the $1{,}031$ concepts. The item bank of $297$ questions and
the unified bank give the same ordering with smaller magnitudes ($1.00 \to 5.48 \to 37.08$ and
$1.00 \to 5.04 \to 74.90$ respectively), so the effect is not a property of one channel.}
\label{tab:dispersion}
\small
\begin{tabular}{lrrr}
\toprule
population & $n$ & distinct $q^\star$ & effective \#questions \\
\midrule
ML-25M cold, $k{=}0$ & 10{,}000 & 1 & 1.00 \\
ML-25M after one standard opener & 10{,}000 & 5 & 2.88 \\
ML-25M after two standard openers & 10{,}000 & 9 & 5.04 \\
Real free-text request, placed & 4{,}431 & 442 & \textbf{95.13} \\
\quad with $\geq 1$ named film & 2{,}639 & 413 & 128.46 \\
\quad text-only (no named film) & 1{,}792 & 171 & 38.23 \\
\bottomrule
\end{tabular}
\end{table}

At the cold start every user has the same state, so the loss-optimal question is one question for
everybody --- a floor that is verified rather than assumed. Two standard openers barely move it:
the population still occupies $112$ distinct states and the optimal next question collapses to an
effective five, with $41\%$ of users sharing a single argmax. A single real request spreads it to
an effective $95$. Adaptive and static policies downstream of a closed opener are choosing between
a handful of distinguishable situations; that is the regime in which Act 1's null was measured, and
it is the regime the field's evaluations inhabit.

\paragraph{The same measurement, read against what the selectors actually did.} Act~1 records that
Thompson sampling gives all $10{,}000$ users a distinct interview by the fourth question, and that
GLM-UCB reaches $5{,}739$ distinct interviews by the sixteenth. Set those against the table above:
downstream of two standard openers there are roughly \emph{five} questions genuinely worth asking.
A selector can therefore manufacture thousands of different interviews while choosing among a
handful of materially different options --- which is what ``divergence without exploitability''
means, stated as two numbers from two independent measurements rather than as an interpretation.

\paragraph{What the placement figure does and does not compare.} The bottom rows of
Table~\ref{tab:dispersion} are a different population from the top rows: $4{,}431$ people who wrote
a free-text request, against $10{,}000$ MovieLens users. The contrast is therefore between
\emph{regimes}, not a within-user effect, and on its own it establishes that placement and
dispersion travel together rather than that one causes the other.

Two guardrails travel with this figure and must not be dropped. First, \textbf{dispersion is not
payoff}: Section~\ref{sec:act1} shows a policy can manufacture ten thousand distinct interviews and
still lose, so a divergent criterion surface is necessary, not sufficient. Second, \textbf{headroom
runs the other way} --- the best \emph{concept} question is worth $+0.0283$ in expected NDCG@10 at
the cold states after two openers and only $+0.0076$ at the request-placed states, because a placed
user is already further along the curve. (Within the concept channel the comparison is like for
like; the request-placed states' best question overall is an \emph{item}, at $+0.0138$, which is
the channel flip of Section~\ref{sec:act4}.) Dispersion and headroom are different quantities and
this chapter never conflates them.

\section{Act 3: the demonstration --- put the reported metric in the split}
\label{sec:act3}

Act 1's null and Act 2's diagnosis together specify one experiment. Take the field's own adaptive
machinery --- Golbandi's ternary interview tree~\cite{golbandi2011adaptive}, which asks a question at
each node and routes users down a LIKE / DISLIKE / DID-NOT-ANSWER branch --- and change exactly one
thing: the scalar by which a candidate question is scored. Everything else is held fixed by
construction, and asserted rather than assumed.

\paragraph{The criterion.} At a node, each of the $1{,}000$ candidates ($500$ popularity-ranked items
and $500$ answerability-ranked concepts, in one argmax) is scored by the node-size-weighted mean of
its three children's cohort-mean \emph{measured} held-out full NDCG@10: fold that child's answers
through the frozen certified recommender and rank the catalogue. Nothing is approximated, no
surrogate appears anywhere, and no ground truth is read at selection time --- the criterion is a
cohort-level quantity on the validation cohort and never touches the evaluated user's labels
(asserted: $0$ of $10{,}000$ users on both cohorts have a held-out target inside the routing lookup).

\paragraph{The tree.} Grown on the $10{,}000$ validation users, then frozen. Depth $5$, $15$ leaves,
$22$ nodes, \texttt{MIN\_LEAF}${}={}200$ as the only stopping rule, and a budget of sixteen charged
questions per user --- the splitters on that user's path plus a greedy tail computed inside their
leaf. Every question is charged and folded: an unanswerable one burns the turn and returns the
tower's unseen-exposure token or the concept channel's refusal token. There is no free side
information anywhere in it, which is what separates this arm from the genre split of
Section~\ref{sec:groupgreedy}.

The root splitter is a \emph{concept}, answered by $74.2\%$ of the cohort and splitting it
$2{,}606 / 4{,}814 / 2{,}580$. The best available \emph{item} split dumps $7{,}045$ of $10{,}000$
users into the did-not-answer branch. That contrast is the answerability result of
Section~\ref{sec:act4}, arriving unbidden inside a different construction.

\paragraph{The comparator, and the gate that makes it one.} The tree is measured against a static
order built by the same greedy machinery, on the same objective, over the same pool, under the same
candidate mask, on the same cohort, and resumed from the tree root's own search heap. Its first pick
is therefore necessarily the tree's root question, and the contract gate is exact: \textbf{at $K{=}1$
the per-user difference is $0.0$ for all $10{,}000$ validation and all $10{,}000$ test users on both
axes}, asserted with equality rather than a tolerance. A tree that ``wins'' at one question, where
conditioning is impossible, is measuring its comparator's selection contract and nothing else.

\begin{table}[t]
\centering
\caption{The NDCG tree against the static order built by the same search on the same objective,
$10{,}000$ held-out test users, full/tail NDCG@10, paired percentile bootstrap ($10{,}000$ resamples,
seed $98765$). $K{=}1$ is exactly zero by construction --- the contract gate, not a result.
Absolute at $K{=}16$: tree $0.26936/0.13074$, static $0.25900/0.12370$, no questions
$0.15855/0.02819$.}
\label{tab:tree}
\small
\begin{tabular}{lcc}
\toprule
$K$ & full gain [95\% CI] & tail gain [95\% CI] \\
\midrule
1 & $0.00000$ (exact) & $0.00000$ (exact) \\
2 & $+0.00528$ $[+0.00324,+0.00725]$ & $\mathbf{-0.00262}$ $[-0.00481,-0.00051]$ \\
4 & $+0.00891$ $[+0.00632,+0.01147]$ & $+0.00482$ $[+0.00214,+0.00756]$ \\
8 & $+0.01062$ $[+0.00775,+0.01354]$ & $+0.00835$ $[+0.00543,+0.01130]$ \\
16 & $+0.01035$ $[+0.00739,+0.01323]$ & $+0.00705$ $[+0.00394,+0.01012]$ \\
\bottomrule
\end{tabular}
\end{table}

\paragraph{The result, with its qualifier attached.} \textbf{Conditioning on cold, closed, charged
questions beats a strong static order on held-out users from four questions onward, on both axes.}
Seven of the eight $K\ge2$ intervals exclude zero in the tree's favour. The eighth excludes zero
\emph{against} it: at two questions the tree is significantly worse on the tail
($-0.00262$, $[-0.00481,-0.00051]$), and the same cell has the same sign on validation without
separating. Conditioning needs a few answers before it pays, and that is part of the finding rather
than a footnote to it.

This is the arm that settles Act 1's ambiguity. A fixed order is an adaptive policy that ignores its
inputs, so no measurement can put adaptivity below staticness --- only particular methods. Here the
particular method is held identical to the published one in every respect except the scalar it
optimises, and the sign of the conditioning gain flips. \textbf{The field's null is about objectives,
not about adaptivity.}

\subsection{The same tree, with six different criteria in its split}
\label{sec:criterion}
\label{sec:huyu}

The mechanism claim above is a claim about one variable, so we vary that one variable. The same tree
code, the same ternary branch rule, the same min-leaf, the same charged behavioural answers, the same
certified recommender, the same candidate mask, the same $1{,}000$-candidate pool, the same two
cohorts and the same sixteen-question budget --- with the split criterion replaced by each of three
published alternatives in turn. Hu and Yu's framework~\cite{hu2013interview} states its split loss
generally and lists a pairwise rank loss among the options; their ranking experiments ran a pointwise
logistic loss; Golbandi's is within-node rating variance. All three are scored through their own
second-order split gain, faithfully, with the gradients and diagonal Hessians checked against central
finite differences before the harness ran. Nothing is sampled: the pairwise loss sums over all
$137{,}591{,}707$ fold-in pairs.

Two shuffle controls complete the set: a uniform draw over the whole pool, and a uniform draw
restricted to candidates whose split clears the minimum leaf size --- the fair control, which holds
the tree machinery fixed and destroys only the choice of question.

\begin{table}[h]
\centering
\caption{One variable changed. Gain against the same static order, $10{,}000$ test users, $K{=}16$,
paired bootstrap. All $56$ non-nested intervals in the six-arm matrix exclude zero.}
\label{tab:splitcrit}
\small
\begin{tabular}{lrrrr}
\toprule
split criterion & full@10 & full gain & tail@10 & tail gain \\
\midrule
measured NDCG@10 (ours) & $0.26936$ & $\mathbf{+0.01035}$ & $0.13074$ & $\mathbf{+0.00705}$ \\
pairwise rank loss~\cite{hu2013interview} & $0.22859$ & $-0.03041$ & $0.10031$ & $-0.02338$ \\
rating variance~\cite{golbandi2011adaptive} & $0.22840$ & $-0.03061$ & $0.10124$ & $-0.02246$ \\
logistic loss (as run)~\cite{hu2013interview} & $0.17274$ & $-0.08627$ & $0.06817$ & $-0.05552$ \\
\midrule
shuffle, min-leaf-respecting & --- & $-0.04817$ & --- & $-0.03412$ \\
shuffle, unrestricted & --- & $-0.03934$ & --- & $-0.02126$ \\
\bottomrule
\end{tabular}
\end{table}

\textbf{Exactly one of the six criteria places the adaptive interview above the static order.}
Changing nothing but the number a question is scored by flips the sign of the conditioning gain. Four
readings follow, and three of them were not what we pre-registered.

\textbf{The separating axis is measured-versus-surrogate, not ranking-versus-rating.} Hu and Yu's
pairwise rank loss and Golbandi's rating variance are \emph{indistinguishable from one another} on
our metric: $+0.00020$ $[-0.00331,+0.00371]$ full and $-0.00092$ $[-0.00448,+0.00277]$ tail in a
paired comparison --- the only null in the whole matrix. Upgrading a rating surrogate to a ranking
surrogate bought nothing measurable. The pre-registered expectation was an ordering; what happened
was $\text{(measured)} \gg \text{(pairwise)} \approx \text{(variance)}$, and the claim narrows
accordingly. Rank-awareness is not the fix. Measuring is.

\textbf{The failure is in the approximation, not in the shape of the loss.} At the root, over all
$1{,}000$ candidates, the \emph{measured} held-out pairwise rank loss --- untruncated, undiscounted,
over the full catalogue --- tracks measured NDCG@10 at Spearman $+0.784$ and forgoes only $+0.00459$,
about $29\%$ of what a random pick forgoes. Hu and Yu's Taylor-approximated \emph{training}-loss
version of the same objective, at the same node, correlates at $-0.149$, ranks the metric's best
question $75$th of $1{,}000$, and forgoes $+0.01473$. Of the ways our criterion differs from theirs,
the one doing almost all the work is that ours is the measured value of \emph{asking} rather than a
training loss over the tree fit. Truncation and the positional discount are real but second-order:
even measured properly, the untruncated loss picks a different first question and recovers $40\%$ of
the metric's top ten.

\textbf{Depth is not virtue.} The min-leaf-respecting shuffle control grows deeper (7 against 5),
wider ($27$ leaves against $15$) and more varied ($286$ distinct questions against $159$) than the
criterion that wins, and still loses by $-0.04817$. No argument from tree size survives that row. The
unrestricted shuffle degenerates to depth $0$ --- its root question is answerable by $4.5\%$ of users
--- and is therefore a randomly chosen static order; the two published criteria beat it on full
($+0.00893$ and $+0.00873$, both separating) and on the tail at sixteen questions neither separates
from it ($-0.00212$ $[-0.00540,+0.00127]$ and $-0.00120$ $[-0.00486,+0.00245]$). Both halves are
reported.

\textbf{The logistic arm is the sharpest single row.} It opens on exactly our question --- the same
concept, the same ternary partition, a $K{=}1$ difference identically zero on all $20{,}000$ users ---
and then degrades monotonically: $-0.00896$ at two questions, $-0.02777$ at four, $-0.05605$ at
eight, $-0.08627$ at sixteen. Its absolute full NDCG@10 peaks at two questions and falls below its own
one-question value by sixteen. Under the loss that paper's ranking experiments actually ran, on
held-out users, with every question charged and every answer folded, asking more questions makes the
recommendation worse. A criterion can agree with the metric on the first, largest, easiest decision
and still walk the interview downhill afterwards.

One structural fact from the same table supports Section~\ref{sec:act4} independently: \textbf{every
criterion that can choose at all splits the root on a concept.}

\paragraph{Two scope sentences that travel with the ablation.} This is a comparison on our ruler. It
establishes that these criteria, transplanted into our harness and given every advantage we give our
own, lose to a static order built by maximising the reported metric; it is not a refutation of the
published results on their own protocols. And it carries a cost asymmetry that runs against us: the
Taylor construction makes gradient and Hessian candidate-independent, so all $1{,}000$ candidates are
scored in a single pass of about eleven seconds, where our exact objective pays a recommender fold
per candidate and cost $37{,}191$ seconds to grow one tree. That is roughly three orders of magnitude
per node; it is a genuine advantage of their construction over ours, and it is exactly why a
parametric policy rather than a tree is the method this line of work should end at.

The ablation is provably on the same ruler as the rest of the chapter: re-measured through the
ablation's own code path, our arm reproduces every published tree number to $2.81\times10^{-7}$.

A faithful snap on Hu and Yu's own protocol (MAP, ten sampled negatives, their top-$k$-popular static
comparator, ML-1M) accompanies the transplant, at depths three and five. Their results are
figure-only, so ``matching'' means the arm ordering, the direction with budget and the gap sizes to
pixel precision, never a published digit. Both readings of their answerability protocol are reported,
because the depth trend depends on which is taken: their training users are about twice as answerable
as the users the method is deployed on, and making the two symmetric moves their
tree-minus-random-questions gap from $+0.00288$ to $+0.00866$ at depth five while shrinking their
tree-minus-fixed-order gap to $+0.0048$.
\PLAN{on their own protocol, four cells remain unrun and are listed as unrun rather than omitted:
depths eight and ten, the pairwise-loss arm of that replication, the RMSE panel with fMF, and
Netflix and Yahoo! Music --- which we do not hold and therefore never claim.}

\subsection{Changing the bank buys more than conditioning does}
\label{sec:bankvsconditioning}

The tree's gain is real and it is small, and the honest way to say how small is to price it against
the other thing one could change. Both sides below are static-versus-static or
adaptive-versus-static on the same cohort, at the same budget, with every order selected under the
scoring contract.

\begin{table}[h]
\centering
\caption{What each change is worth, $10{,}000$ test users, full NDCG@10, against the interview's own
lift over asking nothing. The bank change involves no conditioning of any kind.}
\label{tab:bankdecomp}
\small
\begin{tabular}{lrr}
\toprule
change & gain (full) & share of the static's own lift \\
\midrule
item pool $\to$ item$+$concept pool, equal budget & $\mathbf{+0.01430}$ $[+0.01124,+0.01739]$ & $16.6\%$ \\
tree conditioning (combined pool, $K{=}16$) & $+0.01035$ $[+0.00739,+0.01323]$ & $10.3\%$ \\
genre conditioning (item pool, $K{=}15$) & $+0.00862$ $[+0.00595,+0.01131]$ & $10.2\%$ \\
\bottomrule
\end{tabular}
\end{table}

On the tail the ordering is not close: $+0.02199$ for the bank change against $+0.00705$ for
conditioning. \textbf{If one thing can be changed, change what the system asks about, not how it
decides what to ask.} That is the sentence this chapter would least like to be true, and it is
measured on the same ruler as everything else.

\subsection{What it costs: two fifths of users are worse off, and validation overstates the rest}

\paragraph{Shrinkage.} Against the same comparator, the tree's gain at sixteen questions is
$+0.02794$ full and $+0.01055$ tail on the selection cohort and $+0.01035$ / $+0.00705$ on test ---
$2.70\times$ and $1.50\times$. The asymmetry survives in direction only: the full-metric edge is the
more heavily validation-fitted half. The overfitting signature itself does not survive. On validation
the leaves that gain most are the small ones (Spearman between leaf size and gain $-0.48$,
$p{=}0.071$); on test that correlation is $+0.10$ ($p{=}0.71$).

\paragraph{Negative leaves.} \textbf{Six of the fifteen leaves lose on test, covering $39.3\%$ of the
cohort} ($3{,}933$ users), against none on validation. This is a property of the construction rather
than a fault in it: a leaf inherits questions chosen for its \emph{parent's} cohort, so the static
order is genuinely not inside a leaf's search space and a leaf can lose. Two leaves crossed from
positive to negative when the comparator was rebuilt, which is what a stronger comparator should do.
There is a suggestive pattern --- three of the negative leaves are among the four
lowest-absolute-NDCG leaves, so the tree helps least where users are hardest to serve --- but the
fourth-lowest carries the largest positive gain of any leaf ($+0.0584$), so there is no clean rule
and we do not state one.

\paragraph{Reach, which conditioning does not touch.} The three children of the root's
did-not-answer branch have baselines $0.374 / 0.320 / 0.087$ against a parent at $0.119$. Answering
the question \emph{at all} --- like or dislike, barely mattering which --- is worth roughly $+0.25$,
where the largest question-\emph{choice} effect anywhere in the tree is $+0.048$. And $1{,}785$ users
answer neither of the first two questions, $1{,}438$ of whom cannot answer the third either:
\textbf{about $14\%$ of the cohort is effectively untouched by the interview}, sitting at $0.087$
while everyone else moves. Section~\ref{sec:act5} is about them.

\subsection{Why the tree stopped at depth five: the rule, not the data}
\label{sec:depth}

Realised depth is the first thing a reader will read as a measurement of how far conditioning pays,
and here it is not one. The tree splits three ways and refuses a split if any child holds fewer than
$200$ users, so before any data is consulted a strictly ternary split needs $n\ge600$ and a
fully-answered binary one needs $n\ge400$. Restated, the floor is a \emph{window on the answer rate},
$[400/n,\;1-200/n]$, and it closes at $n=600$. Admissible candidates fall from $876$ of $1{,}000$ at
the root to $0$ at $n=660$.

Classifying all fifteen stopped nodes against the artefacts rather than by inference: \textbf{five
stopped on size at $n<400$} (no split of any shape can clear the floor), \textbf{three on size at
$400\le n<600$}, \textbf{seven on shape} (the node had the users, but the argmax splitter put a child
under the floor), and \textbf{none --- zero --- for want of a positive-gain split}. Every stop had a
strictly positive gain already in hand, from $+0.00672$ to $+0.04251$.

The binding child is also legible. In \textbf{eight of the fifteen stops it is the did-not-answer
branch of a concept that $92.7$--$100\%$ of the node could answer} --- three of them at exactly
$100\%$, so that branch is empty. The other seven are the mirror image: an item whose answer rate is
too \emph{low} ($0.097$--$0.435$), so the like and dislike branches starve while the no-answer branch
swallows the node. The property that makes a concept worth asking is the property that ends the
branch.

That interaction is measurable at the root as well, and its sign depends on the channel, which is
worth stating precisely rather than in the convenient direction. Among the $1{,}000$ root candidates
the excluded ones are systematically \emph{more} answerable than the survivors (median answer rate
$0.986$ against $0.309$; all $72$ excluded concepts are excluded on their did-not-answer branch; the
most answerable surviving concept sits at $0.9797$, exactly the arithmetic ceiling $1-200/10{,}000$).
Within the concept channel that penalty holds at nine of the twelve nodes with a feasible concept.
\textbf{Pooled over both channels below the root it runs the other way} ($+0.117$): the absolute
floor becomes a larger fraction of a smaller node and starts eliminating low-answer-rate items
wholesale. Both are the same window closing from opposite sides, and the answerability penalty is a
\emph{concept-channel} phenomenon, not a pooled one.

\textbf{What this licenses is a negative in both directions.} This run supports no claim that
conditioning stops paying at depth five, and no claim that a deeper tree would pay. The banked
evidence argues against extrapolating: the three depth-five leaves score $+0.028$, $-0.005$ and
$-0.006$ on test.

\subsection{The tree partitions and diverges. It is not shown to route}
\label{sec:routing}

The attractive reading of an adaptive interview is a curriculum --- a broad answer determining which
finer question follows. We tested it, against controls fixed in a design sheet written before any
statistic was computed, and it did not survive.

\textbf{Divergence is decisive.} Across the $105$ leaf pairs, the median overlap of \emph{free} picks
(a leaf's questions minus the inherited path) is $0$: \textbf{eighty of the $105$ pairs share nothing
at all}, the median leaf differs from another in $14$ of its $15$ positions, and $180$ free slots
yield $154$ distinct questions. Two sets drawn at random from the leaves' own question union would
overlap $0.935$; the observed $0.267$ is $3.5\times$ below even that matched chance level. Only nine
of the global static order's fifteen questions appear anywhere in the tree, and four leaves share
none of them. That is what the conditioning gain buys: a question set the static order never reaches.

\textbf{Congruence is a null.} The pre-registered test asks whether the questions below a branch are
\emph{about} the answer that created it, measured as similarity in a $1{,}128$-dimensional tag
relevance space between a node's splitter and the questions asked below each of its children. Pooled
over the seven split nodes, like-branch congruence is $0.7705$ against a matched-availability control
of $0.7691$ $[0.7572,0.7821]$, $p{=}0.82$; the like-minus-dislike contrast is $+0.0054$ against a
branch-label shuffle of $[-0.0200,+0.0200]$, $p{=}0.61$. Only four of the seven nodes have a positive
contrast. Both legs of the decision rule fail. A question drawn at random from the pool, matched on
channel, answer rate and divisiveness, is exactly as related to the splitter as the question the tree
actually chose.

\textbf{Therefore no routing or curriculum claim is made about the tree anywhere in this chapter.}
The adaptive gain is real and its questions are genuinely different; \emph{why} those questions are
the right ones for that branch is not something we can currently characterise. One secondary cell did
separate --- at the root only, on items only, users who liked the opener are subsequently shown items
more relevant to it than a matched-random control ($0.4348$ against $0.3379$ $[0.2768,0.4053]$,
$p{=}0.004$) --- and it is recorded here because suppressing a positive that failed the
pre-registered test would be the same error in the other direction. It is one node of seven, one of
six secondary tests with no multiplicity correction, and it disappears under the primary statistic.
It is a hypothesis, not a result.

\section{Act 4: answerability orders the channels}
\label{sec:act4}

A static order is not a strategy. Nothing in it informs anything: the fifteenth question is fixed
before the first answer arrives, so any language about a curriculum is a category error when applied
to it. What a static order's \emph{composition} can show is which questions a loss-maximising search
reaches for first, and here it reaches for the coarse channel --- for a reason we can measure.

Three independent measurements say the same thing.

\textbf{(i) The greedy order's composition.} The prefix-optimal sequence over the joint bank
(Table~\ref{tab:greedy}, combined arm) opens \texttt{c\,c\,i\,c\,c\,i\,i\,c\,c\,c\,c\,i\,i\,i\,c\,c}:
concepts hold ten of the sixteen slots and both of the first two, and the item block arrives in the
middle rather than at the front. Nobody specified that ordering; the search is a plain argmax over
measured NDCG@10 at each prefix.

\textbf{(ii) The per-state criterion.} In the loss-coupled criterion of Section~\ref{sec:act2},
where both channels are scored in one unit, the best available question at cold states is a concept
by more than a factor of twelve in expected gain ($+0.0238$ against $+0.0019$ for the best item), and
at request-placed states the optimum flips to an item ($+0.0138$ against $+0.0076$ for the best
concept).

\textbf{(iii) The tree's root.} The adaptive tree of Section~\ref{sec:act3}, choosing freely from
both banks, opens on a concept answered by $74.2\%$ of the cohort, where the best item split sends
$70\%$ of users into the did-not-answer branch.

\paragraph{The mechanism is answerability, and it is the same one that explains Act 1's failures.}
Early, a specific question is uninformative \emph{because the cold user cannot answer it}: sixteen
item questions collect $3.11$ answers where sixteen concept questions collect $12.36$, a factor of
four. A broad question is coarse but answerable. This is a statement about which channel a
loss-maximising selector reaches for at which budget, measured on held-out users --- not a claim that
anything is being taught, and not a claim about routing, which Section~\ref{sec:routing} settles
negatively.

\paragraph{Which single channel is the better bank.} At sixteen questions on the test cohort, with
all three orders selected under the scoring contract, \textbf{concepts win the tail decisively and
tie on full}: concepts minus items is $-0.00003$ $[-0.00315,+0.00310]$ on full and $+0.01436$
$[+0.01093,+0.01776]$ on tail, and the combined bank beats both on both axes ($+0.01430$ / $+0.02199$
over items, $+0.01434$ / $+0.00763$ over concepts). The channel ordering depends on budget, and that
dependence is the answerability story rather than a wobble: \textbf{concepts lead on full at every
budget below sixteen} ($+0.00481$, $+0.00781$, $+0.01035$, $+0.01116$ at one, two, four and eight
questions, all excluding zero) and the item order only catches up at sixteen. Cold users can answer
concepts; the item bank needs a longer interview to find questions they have seen. At one and two
questions the combined order \emph{is} the concepts order, so those cells are zero by nesting.

An earlier draft read this comparison the other way --- concepts leading on both axes at sixteen ---
from validation figures under a candidate pool defined by what each arm happened to ask. An item
question that gets answered deletes that item from the candidate list and a concept question deletes
nothing, so that pool flattered whichever arm asked items. Under the protocol pool, on held-out
users, with both orders selected under it, the full axis is a dead heat.

The composition remains visibly steered by the objective: every order here maximises \emph{full}
NDCG@10, and the tail result is a secondary outcome of a search that never optimised it.

\section{Act 5: reach --- the users a closed interview cannot touch}
\label{sec:act5}

Act 3 settles differentiation and leaves reach untouched. About $14\%$ of the cold cohort answers
nothing across sixteen charged questions and sits at $0.087$ full NDCG@10 while everyone else moves,
and inside the tree the value of answering \emph{at all} ($\sim\!+0.25$) dwarfs the value of the best
question \emph{choice} ($+0.048$). No amount of conditioning helps a user who cannot answer: every
selector in this chapter, adaptive or static, is choosing among questions that this population will
refuse. The only instrument that can reach them is a question they can answer in their own words.

That yields a falsifiable intervention: replace the closed first question with an open one, and
re-run the identical baselines. The question distribution is mined from $4{,}431$ real free-text
movie requests (Section~\ref{sec:reddit}); the users, the answers and the held-out scoring remain
ML-25M throughout, so ground truth is never compromised. The primary analysis is a
difference-in-differences: (adaptive $-$ static) under an open opener, minus the same difference
under a closed opener. A validation gate runs first: the matched-initialisation dispersion must
reproduce the request-state dispersion of Table~\ref{tab:dispersion}, or the bridge leaks and the
analysis stops.

\PLAN{the open-opener runs (favourite-film, favourite-genre, and a distinctiveness-framed variant),
the difference-in-differences with paired CIs, and the dispersion validation gate. The primary
outcome is \emph{reach} --- what fraction of the answer-nothing population the open opener converts,
and what that conversion is worth --- not a starved-opener story about the closed selectors, which
Act 3 has already answered. Comparators must include an init-aware-but-fixed static order, the best
fixed order fitted on the initialised population, or we commit the comparator-construction error this
chapter criticises.}

\subsection{The corpus}
\label{sec:reddit}

The open-opener distribution and the dispersion measurement both rest on a corpus of real requests:
$4{,}431$ threads, each an unprompted free-text request with ground-truth items, linked to the
catalogue by a two-tier pipeline (deterministic rules for unambiguous spans, a language model
adjudicating the remainder) whose precision we measured against a $500$-mention hand-labelled gold
standard: $0.572$ for the corpus's own links, $0.960$ for our pipeline. Every request is parsed into
typed preference slots ($94.0\%$ are genuine requests; median five slots; positive anchors in
$63\%$ of threads, explicit dislikes in $16.5\%$). These artefacts serve three roles: they supply
the opener distribution for Act 3, they are the ecological endpoint of the openness gradient in
Table~\ref{tab:dispersion}, and their volunteered slots form the answer pool for a withheld-slot
replay in which the interview must recover what the user would have said. \PLAN{the withheld-slot
replay, both the pool-restricted refusal-free variant and the refusal-counting variant}

\subsection{What people actually volunteer}
\label{sec:opener_census}

The corpus supports a question no ML-25M measurement can answer: \emph{what does a real user say
first?} We parse all $4{,}431$ threads into typed slots and report the distribution. Nothing here is
scored; it is a census, and it constrains what an opener can be.

\begin{table}[t]
\centering
\caption{What $4{,}431$ real free-text requests contain. \emph{Marginal} counts a slot appearing at
all; \emph{opener} counts it being the highest-priority slot present, under the priority order
\texttt{anchor} $>$ \texttt{genre} $>$ \texttt{tone} $>$ \texttt{era} $>$ \texttt{constraint}
declared before the count. The priority ordering is a modelling choice; the marginal is not.}
\label{tab:opener_census}
\footnotesize
\begin{tabular}{lcc cc}
\toprule
slot & \multicolumn{2}{c}{marginal} & \multicolumn{2}{c}{as opener} \\
\cmidrule(lr){2-3}\cmidrule(lr){4-5}
 & threads & share & threads & share \\
\midrule
positive anchor (names a film) & $2{,}793$ & $63.0\%$ & $2{,}793$ & $63.0\%$ \\
genre                          & $1{,}763$ & $39.8\%$ & $651$     & $14.7\%$ \\
constraint                     & $1{,}360$ & $30.7\%$ & $230$     & $5.2\%$ \\
tone                           & $946$     & $21.4\%$ & $146$     & $3.3\%$ \\
negative anchor                & $730$     & $16.5\%$ & ---       & --- \\
era                            & $433$     & $9.8\%$  & $86$      & $1.9\%$ \\
negative genre                 & $184$     & $4.2\%$  & ---       & --- \\
\midrule
\multicolumn{5}{l}{no slot of the five present: $525$ threads ($11.9\%$)} \\
\bottomrule
\end{tabular}
\end{table}

Three findings bear directly on the design of an opener.

\textbf{Requests are bundles, not single facts.} $93.6\%$ of threads carry at least two slots
(median four). Taking the four routable slot types together, $35.9\%$ of threads give both anchor
films \emph{and} a genre, tone or era; $27.1\%$ give films alone; $19.9\%$ give a signature with no
film; $17.0\%$ give neither. A design that assumes the user states one thing is modelling a minority
of the traffic, and the $17.0\%$ who state nothing routable are the population for whom a cold
interview remains the only instrument.

\textbf{What people volunteer first is a film, and a film is evidence rather than a question.} The
single most common opener is a title, in $63.0\%$ of threads. A named film is folded directly by the
recommender; it is not a question the system chose to ask. This is the ecological counterpart of the
result in Section~\ref{sec:act3}: volunteered items pay through the fold. It also bounds what
routing can be keyed on --- $5{,}868$ distinct titles appear across $2{,}793$ threads, far too
granular to define a cohort, so anchors enter as evidence and only genre, tone and era can carry a
route.

\textbf{Most of what is said is expressible in our vocabulary, and the residual is informative.}
Every one of the $7{,}572$ OP-named films ($100.0\%$) lies inside the $18{,}359$-item training
vocabulary, so a request is representable in the action space without approximation. Mapping the
free-text values onto attributes we hold succeeds for $95.1\%$ of genre mentions, $70.2\%$ of eras
and $66.1\%$ of tones. Folding films and mapped concepts together, $3{,}078$ of $4{,}431$ requests
($69.5\%$) can be placed as a state. The unmapped residual is a measured price of the attribute
vocabulary rather than a nuisance: it is concentrated in \emph{coming of age}, \emph{Christmas},
\emph{supernatural} and \emph{dystopian} on the genre side, and on the tone side in precisely the
affect words with no genome tag --- \emph{wholesome}, \emph{uplifting}, \emph{cozy} --- alongside
relative eras (\emph{recent}, \emph{newer}) that denote no fixed decade.

\textbf{What is requested and what is watched are close to inverted.} Horror is $27.5\%$ of mapped
genre mentions, the largest single category, followed by Thriller at $15.8\%$. In the ML-25M
validation cohort, partitioned by the same fold-in genre rule used in
Section~\ref{sec:groupgreedy}, Horror is the favourite genre of $51$ of $10{,}000$ users and Drama of
$5{,}481$. The population that asks and the population we can score are not the same population.
This is a scope condition on any elicitation result measured on ratings data alone, and it is the
reason a benchmark weighted by observed request frequencies is not interchangeable with one weighted
by corpus prevalence.

\PLAN{the slot extraction is model-produced and, unlike the linking, carries no gold-set accuracy
number; a hand-labelled sample of requests against their extracted slots is required before these
shares are load-bearing.}

\subsection{What a user's state actually encodes: the release-era axis}
\label{sec:era}

If an interview is to start from where a user already stands, the states of users who stand in
similar places must be similar for reasons that a user could \emph{say}. That turns out not to be
free.

Folding each of the $140{,}768$ training users' liked history through the frozen tower gives a
state per user, and partitioning that space yields large, stable groups. Their content, however, is
not taste. Decoding each group's centroid through the recommender and reading off the films it
favours \emph{relative to the average user} --- the relative step matters, since absolute scores
return the globally popular list for every group --- gives clusters defined by a window of years:
one of $1995$--$96$ social dramas, one of $2009$--$14$ young-adult adaptations, one of $2013$--$15$
prestige releases. The partition is chronological, not thematic.

The cause is a property of the data rather than of the model. A user joins MovieLens at some point
and rates what is current, so their history encodes \emph{when they were active} at least as
strongly as what they enjoy. Two people with identical taste who arrived a decade apart occupy
different regions of the state space.

The effect is directly measurable. Every catalogue item carries a release year, so each principal
direction of the state space can be decoded to per-item scores and correlated against it. The
second direction reaches $\rho = +0.709$ (Spearman, $14.1\%$ of variance); the first reaches
$+0.362$ ($29.1\%$) and the third $-0.352$ ($6.2\%$). Close to half the variance is era-loaded.
Underneath, the fourth direction correlates at $+0.036$ --- genuine preference structure, carried by
a component roughly a third the size of the chronological one.

Because the year of every item is known, the offending direction can be constructed rather than
guessed. Writing $y$ for the standardised release year across items and $W^{d}$ for the decoder,
$d_{\text{era}} = (W^{d})^{\top} y$ is the state direction whose decoded scores align maximally with
year; measured on its own, it correlates at $\rho = +0.875$, confirming the construction. Projecting
this single direction out of every user state removes almost all of the effect --- the leading
directions fall to $+0.061$, $+0.169$ and $-0.113$ --- while costing $2.5\%$ of state norm
($10.87 \rightarrow 10.60$). It is a thin direction with outsized influence, and a targeted
projection is preferable to discarding whole components, which would throw away the preference
signal the first direction carries alongside its chronological loading.

What is left behind is nameable. The same partitioning applied to the corrected states produces
groups whose most distinctive films and whose most distinctive attribute agree: a superhero group
(\emph{Iron Man 3}, \emph{Avengers: Age of Ultron}, \emph{Captain America}) whose sharpest attribute
in the request vocabulary is \texttt{marvel}; an art-house group (\emph{Holy Motors}, \emph{The White
Ribbon}, \emph{A Field in England}) named \texttt{art house}; a dialogue-led group (\emph{Birdman},
\emph{Whiplash}, \emph{The Grand Budapest Hotel}) named \texttt{dialogue driven}; and further groups
named \texttt{gore}, \texttt{inspiring} and \texttt{teenagers}. In four of the six, the attribute that
separates the group best \emph{within the vocabulary real requesters use} is the same attribute that
separates it best within the full $992$-attribute vocabulary --- the plain word costs nothing.

Two caveats travel with this. A single projection is not a complete correction: several groups
survive the selection criterion while still being defined by a late-1990s window, and they are
reported rather than used. And the correction is measured on the training population only; no
held-out ranking claim rests on it here.

\subsection{Groups a user could name, and where real requests land}
\label{sec:clusters}

Correcting for era leaves a partition worth routing on. Fitting on the $140{,}768$ training users
and sweeping the number of groups, each group is then described in two independent ways: by the
\emph{attribute} that most separates it --- the share of an attribute's endorsers falling inside the
group, minus the largest share falling in any other, estimated on a held-out half --- and by the
\emph{films} it favours relative to the average user. The two descriptions are produced from the
same states but by different routes, so their agreement is a check rather than a restatement.

\begin{table}[t]
\centering
\caption{Six groups at $k{=}20$ whose attribute and films agree, fitted on train. \emph{Plain word}
restricts the attribute vocabulary to the $341$ words that appear in real requests
(Section~\ref{sec:opener_census}); ``same'' means the unrestricted and restricted choices coincide.
\emph{Requests} is the share of the $3{,}078$ foldable real requests whose own state lands nearest
this group. The final column is that share divided by the group's share of training users.}
\label{tab:clusters}
\footnotesize
\setlength{\tabcolsep}{4pt}
\begin{tabular}{llll rr r}
\toprule
attribute & plain word & most distinctive films & $n$ & train & requests & ratio \\
\midrule
\texttt{inspiring} & same & Dark Knight Rises, Interstellar, Django & $9{,}743$ & $6.9\%$ & $\mathbf{20.2\%}$ & $2.9\times$ \\
\texttt{dialogue driven} & same & Birdman, Whiplash, Grand Budapest Hotel & $7{,}475$ & $5.3\%$ & $\mathbf{16.8\%}$ & $3.2\times$ \\
\texttt{gore} & same & Dredd, Watchmen, The Hateful Eight & $5{,}825$ & $4.1\%$ & $\mathbf{16.0\%}$ & $3.9\times$ \\
\texttt{arms dealer} & \texttt{marvel} & Iron Man 3, Avengers: Age of Ultron & $6{,}882$ & $4.9\%$ & $11.5\%$ & $2.4\times$ \\
\texttt{teenagers} & \texttt{school} & Aquamarine, Princess Diaries 2, Frozen & $4{,}945$ & $3.5\%$ & $6.1\%$ & $1.7\times$ \\
\texttt{art house} & same & Holy Motors, The White Ribbon, Holy Motors & $4{,}576$ & $3.3\%$ & $5.1\%$ & $1.6\times$ \\
\midrule
\multicolumn{4}{l}{\emph{six groups, combined}} & $28.0\%$ & $\mathbf{75.7\%}$ & $2.7\times$ \\
\bottomrule
\end{tabular}
\end{table}

Table~\ref{tab:clusters} reports the six groups whose two descriptions agree. In four of the six the
separating attribute is already a word real requesters use, so restricting the vocabulary costs
nothing; in the other two the plain substitute is close (\texttt{marvel} for \texttt{arms dealer},
\texttt{school} for \texttt{teenagers}) and loses little margin.

The final columns are the part that could have gone wrong and did not. Group membership was
determined entirely from MovieLens histories; nothing in the procedure saw a request. Yet folding
each real request's own volunteered bundle and asking which group its state falls nearest,
\textbf{these six groups receive $75.7\%$ of real requests while containing $28.0\%$ of users} --- a
concentration of $2.7\times$, rising to $3.9\times$ for the \texttt{gore} group. The groups we
identified as coherent are the groups real people arrive in.

The converse holds equally. Four further groups clear the selection criterion but fail the agreement
check, their distinctive films spanning $1997$--$2000$ with attributes to match (\texttt{1930s},
\texttt{70mm}): residual chronology the single projection did not remove. \textbf{They receive
between $0.7\%$ and $1.3\%$ of requests each}, and the largest training group of all --- $21{,}477$
users, $15.3\%$ of the population --- receives $1.5\%$. Chronological groups are exactly the ones
nobody asks from.

Partitioning the request states themselves rather than assigning them gives an adjusted Rand index of
$+0.269$ against that assignment: agreement well above chance and far from complete. Five of eight
request groups map onto a single taste group with matching content --- horror titles onto
\texttt{gore}, superhero titles onto \texttt{marvel}, art-horror onto \texttt{dialogue driven} ---
while three are diffuse, and all three of those are again defined by late-$1990$s films. An opener
places a user approximately, not exactly, and the gap between $+0.269$ and agreement is the room an
adaptive interview has to work in.

Two groups deserve their failure recorded. One is describable but not separable: its films
(\emph{Cinderella}, \emph{27 Dresses}, \emph{Sisterhood of the Traveling Pants 2}) are plainly teen
romance and its best attributes (\texttt{fairy tale}, \texttt{whimsical}, \texttt{teen}) clear the
purity requirement, but every one carries a \emph{negative} margin because the \texttt{teenagers}
group already owns those words --- one taste region split in two, and correctly merged. The other
cannot be described at all: its best attributes (\texttt{goofy}, \texttt{silly}) fall below the
purity requirement with margins from $-0.077$ to $-0.115$, and its films are unrelated obscurities.
It is a group of people who rate unpopular things, which is a behaviour rather than a preference,
and no question a user could answer would route anyone to it.

\subsection{What volunteered information is worth}
\label{sec:volunteered}

Everything so far establishes that people volunteer, that what they volunteer places them, and that
the places are nameable. This section asks the operational question: \emph{what is it worth?} Three
measurements, on three different objects, answer it --- and they do not agree in the way one might
expect.

\paragraph{It pays through the recommender, not the interview.}
Seeding a cold user with $k$ items drawn from their own history, then running the \emph{identical}
sixteen questions, is worth $+0.0206$ full and $+0.0353$ tail NDCG@10 at $k{=}8$
(Section~\ref{sec:act3}). The questions do not change; the ranking does. A single stated attribute is
worth more still: at the entry state of the \texttt{art house} group, folding the attribute and
asking nothing is worth $+0.056$ over the same cohort's cold start --- larger, on its own, than the
entire fifteen-question conditioning gain the tree earns over a static order.

\paragraph{It pays a second time through the questions --- but only where the fact is specific.}
The first channel is the recommender consuming evidence. The second is the interview \emph{changing
because} of it, and isolating that requires folding the volunteered attribute into \emph{both} arms
so that the evidence itself cancels. Table~\ref{tab:volunteered} does exactly that: it compares each
group's dedicated interview against the cold interview, at matched budgets, with the attribute folded
on both sides.

\begin{table}[t]
\centering
\caption{The second channel: a group's dedicated interview minus the cold interview, at matched
budgets, with the volunteered attribute folded into \emph{both} arms so its own contribution cancels.
Full NDCG@10; $\star$ marks a paired bootstrap interval excluding zero ($10{,}000$ resamples).
Validation users; both arms are fitted on this cohort, so levels are optimistic and the comparison,
not the level, is the measurement.}
\label{tab:volunteered}
\footnotesize
\begin{tabular}{lr rrrrr r}
\toprule
volunteered & $n$ & $K{=}1$ & $K{=}2$ & $K{=}4$ & $K{=}8$ & $K{=}15$ & tail@15 \\
\midrule
\texttt{dialogue driven} & 730 & $+.0044$ & $+.0134^\star$ & $+.0170^\star$ & $+.0137^\star$ & $\mathbf{+.0163^\star}$ & $-.0086$ \\
\texttt{teenagers} & 786 & $+.0123^\star$ & $+.0167^\star$ & $+.0084$ & $+.0154^\star$ & $\mathbf{+.0143^\star}$ & $+.0234$ \\
\texttt{marvel} & 734 & $+.0219^\star$ & $+.0125^\star$ & $+.0133^\star$ & $+.0077$ & $\mathbf{+.0106^\star}$ & $+.0055$ \\
\texttt{art house} & 587 & $+.0043$ & $+.0052$ & $+.0111^\star$ & $+.0098^\star$ & $+.0058$ & $+.0107$ \\
\texttt{gore} & 747 & $+.0028$ & $+.0040$ & $+.0095^\star$ & $+.0114^\star$ & $+.0056$ & $-.0141$ \\
\texttt{inspiring} & $1{,}822$ & $+.0087^\star$ & $+.0012$ & $-.0004$ & $+.0003$ & $+.0051$ & $+.0014$ \\
\bottomrule
\end{tabular}
\end{table}

Seventeen of the thirty cells separate from zero and \textbf{every one of them is positive}; none
separates negative, and $29$ of $30$ point estimates are positive. So the second channel is real: a
volunteered attribute changes which questions are worth asking, and the change is worth up to
$+0.016$ --- roughly forty per cent again on top of the $+0.0399$ that conditioning earns over a
static order.

\paragraph{What predicts the second channel is specificity, not sophistication.}
The pattern across Table~\ref{tab:volunteered} is not the one we expected. \texttt{inspiring} has the
largest cohort and by far the most elaborate interview --- four levels of branching, twenty-one
nodes --- and it separates only at $K{=}1$, adding nothing measurable thereafter. \texttt{marvel},
whose interview branches once, separates at four budgets of five. The property that predicts payoff
is how narrowly the volunteered fact locates a person, not how much structure the system builds on
top of it. A broad attribute buys a better opening question and little else.

\paragraph{Two groups win at the head and lose at the tail.}
\texttt{dialogue driven} ($-.0086$) and \texttt{gore} ($-.0141$) improve full NDCG while losing
slightly on tail, and this must be said plainly rather than reported as a full-NDCG headline: for
those groups the conditioned interview is finding popular relevant items faster and long-tail items
marginally less well. \texttt{teenagers} is the counterexample, with its largest gain on the tail
($+.0234$).

\paragraph{Scope, and a bias that runs in our favour.} Table~\ref{tab:volunteered} is measured on the
users both arms were fitted on, and the two arms are \emph{not} equally advantaged. Each group's
interview was fitted on precisely the cohort it is scored on; the cold interview was fitted on the
full ten thousand and merely happens to cover that cohort. The specialised arm therefore enjoys an
advantage the incumbent does not, and the comparison is biased towards it. These numbers are an
\emph{upper bound} on what conditioning is worth, not an estimate of it.

Two things keep the table informative despite that. The bias is not decisive: in an earlier
experiment, orders fitted per genre on their own cohorts nonetheless \emph{lost} to the unconditioned
tree at those same states, so being specialised does not by itself produce a win. And the direction
of the remaining uncertainty is known --- a tree fitted on the cohort it is scored on lost roughly two
thirds of its measured advantage when carried to held-out users
(Section~\ref{sec:act3}), so the honest expectation is that these effects shrink materially rather
than reverse. \PLAN{the held-out confirmation on the reserved test cohort, which is what converts
this table from an upper bound into a measurement; every number here is provisional until then.}

\PLAN{the screen --- whether, at these groups' entry states, any question improves on the tree's own
first question by more than the pre-registered margin. Note the power constraint: a group holds
$3$--$7\%$ of the $10{,}000$ validation users, so a paired interval at group scale is roughly four
times wider than at cohort scale, and the screen is reported as a selection step validated on a
held-out half rather than as a significance claim at $+0.004$.}

\section{Limitations}
\label{sec:limits}

\textbf{The criterion in Section~\ref{sec:act2} scores against the recommender's own predictive
distribution}, so it measures self-consistent expected gain: a divergence in what the model thinks
is worth asking is necessary but not sufficient for a divergence in what pays. A ground-truth-coupled
criterion would be clairvoyant and could not be computed at interview time.
\textbf{Whether a question is answerable is inferred from the same history that scores the user}
--- we treat ``rated it'' as ``could have answered about it''. Answer \emph{values} are behavioural
throughout (Section~\ref{sec:setup}); it is the answer\emph{ability} judgement that is a modelling
assumption, it is the standard protocol in this literature, and since answerability is load-bearing
for two of this chapter's mechanisms it is stated here rather than left implicit.
\textbf{The adaptive selectors maintain their own posteriors}, because the certified tower is
deterministic and has none --- so no selector here can use the instrument's own uncertainty. That
is a real gap in the interface and it is where a continuous-action policy would live; it is out of
scope for this chapter by design.
\textbf{One free constant} (the UCB exploration weight) was fixed at one posterior standard
deviation and never tuned; tuning it on validation users is open work.
\textbf{The written-request corpus bounds open-answer richness from above}, since a composed post
is not a live spoken answer, while the favourite-film opener bounds it from below; the claim ceiling
is a population conditioned on request-shaped evidence, never ``for real requesters''.

\section{Positioning}
\label{sec:positioning}

\subsection{What we do not claim, stated before anything we do}

Three sentences a reader of this literature might expect are false, and we delete them rather than
hedge them.

\textbf{We are not first to put a ranking objective inside an interview tree's split.} Hu and
Yu~\cite{hu2013interview} state a general split loss that admits a pairwise rank loss, and The
FacT~\cite{tao2019fact} instantiates a BPR term in one --- explicitly framed as an interview, with a
third ``unknown'' child, and carried into a conversational tree by its
successor~\cite{haque2022factcrs}. Both compute their ranking term on observed training interactions,
undiscounted and untruncated.

\textbf{We are not first to use a truncated, position-discounted ranking metric as a selection
criterion.} Bilgic and Bennett~\cite{bilgic2012activequery} instantiate DCG@10 and NDCG@10 inside an
acquisition function; ELO-DCG~\cite{long2010elodcg} has the positional discount without the
truncation, and states our thesis in IR eight years before
Sepliarskaia~\cite{sepliarskaia2018spq}. Related work selects by estimated loss reduction under the
ranker's own training objective~\cite{donmez2008rankal} or selects judgments to estimate a
metric~\cite{carterette2006minimal}. Greedy forward selection with an NDCG wrapper is, separately,
ordinary practice in learning-to-rank feature selection. The technique is not the contribution.

\textbf{And we are not first to test adaptivity against a metric-optimised static order.} Golbandi,
Koren and Lempel did exactly that in 2011: their ternary interview tree~\cite{golbandi2011adaptive}
beat GreedyExtend~\cite{golbandi2010bootstrapping}, a static seed set built by greedily minimising
RMSE with the recommender refit on seed-restricted data --- the reported metric of its era --- and
beat it by a wide margin, six questions against more than twenty. \textbf{That result is evidence
\emph{for} this chapter's diagnosis, not against it.} In 2010--11 the split criterion (within-node
squared loss) essentially \emph{was} the reported metric, so the adaptive method and its static
comparator optimised the same thing and the comparison was fair. The field then moved its reported
metric to truncated top-$k$ ranking and kept the rating-era criteria, and nobody rebuilt the static
comparator on the new metric. So the surviving claim is scoped: \emph{since the reported metric moved
to truncated, position-discounted top-$k$ ranking, no adaptive selector has been tested against a
static order built by maximising that metric.}

What is unoccupied is the conjunction: an adaptive interview whose split criterion is the
\emph{measured} value of a question asked of a cold user, scored as cohort-mean held-out NDCG@10 over
the full catalogue with nothing sampled, with unanswerable questions charged, grown on one cohort,
reported on a disjoint one, and tested against a static order built by maximising that same quantity.
Every clause in that sentence is load-bearing and each is met by a different neighbour.

\subsection{The neighbouring families}

\paragraph{Rating elicitation --- our own line, and the strongest prior art.} Interview
trees~\cite{golbandi2011adaptive,hu2013interview,sun2013learning}, functional matrix
factorisation~\cite{zhou2011functional}, seed-set and volume
optimisation~\cite{liu2011wisdom,fonarev2016rectangular,anava2015budget}, and the recent
learned-static and region-based successors~\cite{kweon2020deep,nguyen2024pere,parapar2021dpe} share
our contract: a cold user, a question budget, and a held-out metric afterwards. All are
surrogate-objective, and we run them rather than citing them. Two deserve individual sentences.
Sepliarskaia et al.~\cite{sepliarskaia2018spq} set out to derive the selection rule from the loss the
system is judged on and \emph{proved} the answer is A-optimality --- because the loss they chose, a
flat pairwise squared error, is a quadratic form in the parameter error. Row five of our own
objective audit is their rule, and we say so rather than double-counting it. The sharpened statement
of what we add follows directly: that equivalence holds for a flat pairwise loss and breaks for a
truncated, rank-discounted one. The positional discount is not decoration; it is what makes the
objective irreducible to a variance criterion. And their paper contains, in 2018, the result Act 1
replicates: static questionnaires beating a dynamic one, explained by the dynamic method optimising
``a function that is different from the loss function, namely weighted generalized variance''. Our
Act 1 is best read as a replication and generalisation of that, not as a discovery.

Greedy SLIM~\cite{proissl2024greedyslim} is the closest protocol sibling that exists --- ML-25M,
no item sampling, NDCG@10, the same short-head tail definition --- and it misses on exactly the axis
we test: its selection objective is a SLIM squared loss, with NDCG@10 reported and never optimised.
Its adaptive foil, a bandit questionnaire over PureSVD~\cite{christakopoulou2016towards}, has NDCG
that \emph{declines} past ten questions on both axes, an independent external corroboration of the
bank-quality effect we measure. Its authors also record the reproducibility hole this battery fills,
and its limitations paragraph identifies unanswerability as the threat to its own offline protocol
without ever pricing it.

\paragraph{Conversational bandits --- the family whose policies genuinely port.} ConUCB and its
successors~\cite{zhang2020conucb,wang2023conlinucb,liu2025clisk} ask about a key-term \emph{or} an
item in one algorithm, with an explicit model of how attribute feedback moves the item-level
estimate, no target item, and regret bounds. That is our item-versus-concept channel, formalised in
2020; CLiME runs on MovieLens-25M with genres as key-terms, so this is not a 2011 straw man. We ran
them (Table~\ref{tab:battery}). What separates them is the quantity: cumulative regret over $T$
rounds against the best arm, on evaluation scales of a few thousand arms, against post-interview
held-out ranking here. Regret prices the probes; we price what the system does after the probing
stops. Neither is wrong and a number in one is not orderable against a number in the other --- so the
acceptance gate for a ported policy is not ``reproduce their published number'' but (i) reproduce the
selection rule as written, (ii) beat random and popularity on our metric, and (iii) label the row
\emph{policy ported, protocol replaced}. One terminological trap is worth naming because it is
theirs, not ours: in parts of this line ``cold start'' means early in the stream, not a user with
zero history.

\paragraph{Deep-RL attribute-asking CRS --- adaptive, and not comparable here.} CRM~\cite{sun2018conversational},
EAR~\cite{lei2020estimation}, SCPR~\cite{lei2020interactive}, UNICORN~\cite{deng2021unicorn},
ConTS~\cite{li2021seamlessly} and the dense 2022--26 successor
layer~\cite{qian2023hutcrs,zheng2025hypercrs} deserve one concession made generously: \textbf{this
family optimises the metric it reports.} Its RL reward is target-hitting and its metric is success
rate, so on the axis this chapter is about it is not guilty of the mismatch we diagnose. Four things
separate it from our contract. First, the simulator answers from the ground-truth \emph{target
item's} attribute set --- a circularity that is \textbf{label-based, not representational}, and
milder in kind than geometric self-play, but with the consequence that every question is answerable
by construction and an unanswerable one costs exactly zero. The quantity Act 5 measures cannot exist
inside that protocol. Second, the conversation is initialised with an attribute \emph{of the target
item} and the candidate set is the items carrying it, so the reported top-ten is a re-rank inside a
pool the answer key created --- and the pool size is not reported anywhere we could find, which is
itself the finding. Third, CRM, EAR, SCPR and ConTS report \textbf{no rank-sensitive metric at all};
hDCG enters with UNICORN, which explains that with one target item per session it is the metric
available. No paper in this line reports a post-interview held-out ranking metric over a full
unsampled catalogue with no known target item. Fourth, the action space and candidate set are both
pruned by construction, and the benchmark is at its ceiling (success rates of $0.99$ at fewer than
six turns), which is why the successor layer competes on turns rather than on success.

Two further observations belong here because they are the family's own. A method with \emph{no RL at
all}~\cite{zhang2022minicorn} --- Monte-Carlo-dropout uncertainty over attribute beliefs and one
hand-written ask-versus-recommend rule --- reports beating SCPR by $67\%$ in success rate, which
makes the contribution of the policy machinery over a good acquisition rule unestablished within the
family's own protocol. And the leakage critique is published from inside the wider
community~\cite{kim2024targetfree}: even attribute-only priming can shortcut recommendation by
implicitly narrowing the candidate space. The honest sentence is not that this evaluation is invalid
but that it prices a different quantity --- and that putting these systems on our contract would mean
replacing the answerer, the candidate pool, the metric and the reward, at which point the reward that
trained the policy no longer exists. We cite this family; we do not run it, and we say why.

\paragraph{LLM elicitation.} The good systems in this line put the language model \emph{outside} the
selection loop --- PEBOL's acquisition is a closed-form Beta-posterior bandit
rule~\cite{austin2024bayesian}, with the LLM extracting an aspect and phrasing the query. What
separates them is scale and contract: PEBOL's catalogue is $100$ items, a cap that exists only for
comparability with an LLM baseline that must fit every item description in a context window, against
our $18{,}359$. There is no benchmark in this family --- $100$ items, $1{,}877$ documents, $30$
articles, or no catalogue at all --- so a reviewer cannot demand the standard one; what a reviewer
can demand is that we justify our protocol, and the contract argument is one we are placed to win.

\paragraph{Meta-learning cold start --- a category error, stated as such.} MeLU~\cite{lee2019melu},
TaNP~\cite{lin2021tanp} and their successors are answer \emph{consumers}, not selectors: the support
set is given, and nothing is asked. Exactly one of them selects, with a single static global ranking,
and in doing so independently reproduced our answerability collapse --- only a quarter of its
evidence candidates were selected by users, against half of the popularity-based ones. It reported
that as an aside and shipped; we cite it beside our own answer rates.

\subsection{Where a neighbour's result runs against ours}

Three, and all three belong here rather than in a limitations paragraph. \textbf{(1)} Golbandi's 2011
win over a metric-optimised static order, reconciled above. \textbf{(2)} Greedy
SLIM's~\cite{proissl2024greedyslim} crossover runs the other way at very short budgets --- their
adaptive bandit wins at five questions and their static questionnaire wins from ten --- on our
dataset with our tail definition. Our tree is exactly zero at one question by construction and
separates from two to four onward, so the shapes are compatible; but we do not write ``adaptive loses
short and wins long'' as though it were a law. \textbf{(3)} A 2025 attribute-aware interview
tree~\cite{gharahighehi2025pairwise} finds genres winning early and items winning from the thirteenth
question --- directionally our channel result, on a different dataset, in RMSE space, attributed to
clustering efficiency rather than answerability. Our differentiator has to be the mechanism and the
ruler, not the direction of the effect.

\subsection{The residual exposure}

\textbf{The comparator is our own construction and no published number validates it.} Nobody has
published a static order built by greedily maximising cohort-mean held-out NDCG@10, so there is no
external number to snap it against, and the usual gate does not apply. Three mitigations, all
checkable: it beats every published selector we ran, against all three rebuilt static orders; the
recipe is deterministic and fully specified, so anyone can rebuild it; and we measured what
comparator weakness costs and reported it against ourselves --- rebuilding the comparator under the
scoring contract moved the argmax at every one of sixteen positions, cut the conditioning gain from
$+0.01371$ to $+0.01035$ on full and from $+0.02092$ to $+0.00705$ on tail, and retracted a
separated interval we had printed in the tree's favour (Section~\ref{sec:act1result}).

\textbf{The effect is small.} The conditioning gain is about a tenth of the interview's own lift and
smaller than changing the question bank with no conditioning at all. \textbf{The validation-to-test
shrinkage is real and asymmetric} ($2.70\times$ full, $1.50\times$ tail). \textbf{And $39.3\%$ of the
test cohort is in a leaf that loses.} None of these is a footnote to the claim; they are the size of
it.

\PLAN{Two named comparators remain to be run on our ruler: the matrix-factorisation seed-set rule
RBMF~\cite{liu2011wisdom} --- the learned static family's other member, DRE~\cite{kweon2020deep},
is built and in Table~\ref{tab:battery} --- and a region-based adaptive method over a frozen
embedding space~\cite{nguyen2024pere}. One further row is designed and worth having as a labelled ceiling:
UpsRec~\cite{lin2023upsrec}, the only published rule that names NDCG as its selection objective,
which reads the evaluated user's own labels at ask time and so can only be run as a declared
privileged oracle.}

\section*{Reproducibility}

Every measurement in this chapter runs on frozen artefacts --- the certified tower \texttt{i26c
ep31} and concept channel \texttt{cv3 ep08} --- under the canonical split, the NO-FLOOR contract,
and behavioural answers; nothing is trained here. Three scope statements belong with that sentence
rather than after it. \emph{Cohorts:} the ML-25M results use all $10{,}000$ test users with no
sampling, and wherever anything is \emph{selected} --- the greedy sequences of
Table~\ref{tab:greedy} --- the selection ran on the $10{,}000$ validation users, with the test
cohort constructed only after the sequences were frozen. \emph{Consumers:} the certified pair is
the recommender under test, but Tables~\ref{tab:act1} and the reversal of Section~\ref{sec:flip}
deliberately score the same interviews with four further recommenders (RecVAE, TaNP, EDDI's
Partial-VAE, Golbandi's leaf model), which is the point of those sections. \emph{Corpus:} the
request-based rows of Table~\ref{tab:dispersion} and all of Section~\ref{sec:reddit} are over
$4{,}431$ linked threads, a different population with its own stated limits. Run records: the
greedy sequences and their intercept gate; the adaptive-versus-static table with its
exact-equivalence regression gate; the tree's contract gate (a $K{=}1$ difference of exactly zero on
all $20{,}000$ users), its reuse gate (per-user vectors and routed row ids asserted bit-identical to
the banked run) and its banked-headline assert; the split-criterion ablation's finite-difference
checks and its two shuffle controls; the dispersion measurement with its shuffle, existential and
determinism controls; the linking pipeline's gold-standard precision; the slot parse.

% ---------------------------------------------------------------------
\bibliographystyle{plain}
\bibliography{../references,../chapterA_v2/chapterA_v2_new_refs}

\end{document}
